% ====================================================================
% graphite_ica_ch1_appendixB_fitting.tex
%   부록 B — 유도식의 실데이터 피팅 워크플로 (구 Chapter 6)
%   흑연 음극 ICA(dQ/dV) 이론서의 보조(supplement) 문서:
%     식별성 · index-1 DAE 수치 해법 · g-grid closure · Plan A/B ·
%     순차 제약(저속 OCV→GITT→다온도 Arrhenius→C-rate) · 반증 워크플로.
%
%   clean-spine 표준(BP plan 2026-06-03)으로 재서술:
%     · G-motive : 새 양/기호는 공식 *앞에서* 물리 말로 동기 부여.
%     · G-interpret: 구한 양은 공식 *뒤에서* 의미·용도(피팅에서 무엇)로 마무리.
%     · G-spine  : BOUNDED/FLAGGED/주의/rigor 단서는 box, 본문 줄기는 박스 없이도 연속.
%     · G-nofwd  : 본문 줄기 forward reference 0(쓰기 전 정의).
%     · G-bridge : 절→절 기계 도약 시 다리 문장; 무거운 수치 형식은 box.
%   물리·수치 내용(S1~S6, index-1 DAE, g-grid closure, Plan A/B, 식별성,
%   AL-60~69) 및 식별자·라벨 네임스페이스(sec:ch6_*, eq:ch6_*)·변수명·한글 보존.
%
% ★★★ 이론장 cross-ref 재동기 필요(BP-series): 본 보조 문서가 참조하는
%   Chapter 1 이론장 식(전하 보존·ξ_eq·k_j·완화 길이 L·spectrum A_L·Θ·kernel·
%   Volterra·피팅 닫힌식)은 새 clean 이론장(graphite_ica_ch1_clean.tex)이 미완이라
%   하드 eq-번호 대신 *의미로* 참조한다. clean 이론장 확정 후 의미 참조 위치를
%   실제 eq 라벨로 재동기해야 한다. (현 standalone 빌드의 undefined ref/cite = 0
%   이어야 하며, 이론장 참조가 의미 참조라 그 조건을 만족한다.)
%
% Date: 2026-06-03  Author: Project_Anode_Fit (Claude 측, 부록 B 보조 트랙)
% ====================================================================
\documentclass[11pt,a4paper]{article}
\usepackage[margin=22mm]{geometry}
\usepackage{kotex}
\usepackage{amsmath,amssymb,mathtools,bm}
\usepackage{booktabs,longtable,array}
\usepackage{enumitem}
\usepackage{amsthm}
\usepackage{hyperref}
\usepackage{fancyhdr}
\usepackage{lastpage}
\usepackage{setspace}
\setstretch{1.13}
\setlength{\parskip}{0.45em}\setlength{\parindent}{0pt}\setlength{\headheight}{14pt}
\hypersetup{colorlinks=true, linkcolor=blue, urlcolor=blue, citecolor=blue,
  pdftitle={흑연 음극 ICA tail — 부록 B (피팅 워크플로, 보조 문서)}}
\pagestyle{fancy}\fancyhf{}
\lhead{흑연 음극 ICA tail — 부록 B (실데이터 피팅 워크플로)}\rhead{\thepage/\pageref{LastPage}}
\renewcommand{\headrulewidth}{0.4pt}

\newtheorem*{groundingbox}{Grounding}
\newtheorem*{boundbox}{유효범위(Validity bound)}
\newtheorem*{flagbox}{FLAGGED (미확정)}
\newtheorem*{keybox}{Keystone}
\newtheorem*{linkbox}{이론장(Ch1--5) 전달식}
\newtheorem*{verifybox}{유도 검증 (차원·극한·부호)}
\newtheorem*{practicebox}{실무 팁 (본문 물리식 아님)}
\newtheorem*{loopbox}{도입 관찰로의 폐합}

\newcommand{\dd}{\mathrm{d}}
\newcommand{\eff}{\mathrm{eff}}
\newcommand{\eq}{\mathrm{eq}}
\newcommand{\app}{\mathrm{app}}
\newcommand{\drive}{\mathrm{drive}}
\newcommand{\bg}{\mathrm{bg}}
\newcommand{\cell}{\mathrm{cell}}
\newcommand{\ext}{\mathrm{ext}}
\newcommand{\OCV}{\mathrm{OCV}}
\newcommand{\tot}{\mathrm{tot}}
\newcommand{\irr}{\mathrm{irr}}
\newcommand{\rev}{\mathrm{rev}}
\newcommand{\ct}{\mathrm{ct}}
\newcommand{\transp}{\mathrm{transport}}
\newcommand{\cb}{\mathrm{cb}}
\newcommand{\dyn}{\mathrm{dyn}}
\newcommand{\sstat}{\mathrm{ss}}
\newcommand{\conv}{\mathrm{conv}}
\newcommand{\stt}{\mathrm{state}}
\newcommand{\hys}{\mathrm{hys}}
\newcommand{\tar}{\mathrm{tar}}
\newcommand{\loss}{\mathrm{loss}}
\newcommand{\obs}{\mathrm{obs}}
\newcommand{\tol}{\mathrm{tol}}
\newcommand{\noise}{\mathrm{noise}}
\newcommand{\brk}{\discretionary{}{}{}}% 표/참고문헌 긴 DOI 줄바꿈 허용(zero-width)
\newcommand{\AL}[1]{\textsf{[AL-#1]}}

\title{\textbf{리튬이온전지 흑연 음극 ICA($dQ/dV$) tail 이론 — 부록 B}\\[0.4em]
\large 유도식의 실데이터 피팅 워크플로 (구 Chapter 6)\\[0.2em]
\normalsize 식별성 $\cdot$ index-1 DAE 수치 해법 $\cdot$ $g$-grid closure $\cdot$ 순차 제약 $\cdot$ 반증 워크플로}
\author{Project\_Anode\_Fit (부록 B 보조 문서, clean-spine 재작성)}
\date{2026-06-03}

\begin{document}
\maketitle

% ====================================================================
% 재동기 flag (문서 상단 1줄, 사용자 지시)
% ====================================================================
\begin{flagbox}
\textbf{★ 이론장 cross-ref 재동기 필요.} 본 보조 문서는 새 clean 이론장
(\texttt{graphite\_ica\_ch1\_clean.tex})이 아직 미완이므로, 참조하는 Chapter 1 이론장 식(전하 보존,
$\xi_{\eq}$, $k_j$, 완화 길이 $L$, spectrum $A_L$, $\Theta$, kernel, Volterra, 피팅 닫힌식)을 \emph{하드
eq-번호 대신 의미로} 참조한다. clean 이론장 확정 후, 본문의 ``이론장에서 유도한 $\cdots$'' 의미 참조를 실제
eq 라벨로 재동기해야 한다. 본 부록 \emph{내부} 식 라벨(\texttt{eq:ch6\_*})과 절 라벨(\texttt{sec:ch6\_*})은
원본 네임스페이스 그대로 보존한다.
\end{flagbox}

% ====================================================================
\section*{\normalfont\Large 부록 B — 유도식의 실데이터 피팅 워크플로 (구 Chapter 6)}
\label{sec:ch1_fitting_practice}
\addcontentsline{toc}{section}{부록 B — 실데이터 피팅 워크플로}
\noindent\emph{본 보조 문서는 이론장(Chapter 1--5)이 확정한 유도식 위에서, 실데이터 피팅을 운영하는
식별성$\cdot$수치 해법(index-1 DAE)$\cdot$순차 제약(저속 OCV$\to$GITT$\to$다온도 Arrhenius$\to$C-rate)$\cdot$
반증 워크플로를 담는다. 즉 ``유도식이 무엇을 말하는가''는 이론장이, ``그 식을 실데이터에 어떻게 거는가''는
본 부록이 맡는다. 가정 원장 AL-60$\sim$69 를 포함한다. 발열$\cdot$히스테리시스 피팅은 각각 이론장 Ch.2$\cdot$Ch.4,
Ch.5 의 정의식을 참조한다.}

% ====================================================================
\section{부록 B 도입 — 이론장 식과 피팅 파라미터 정리}\label{sec:ch6_inherit}
% ====================================================================
\textbf{질문.} 피팅을 시작하기 전에 먼저 답해야 할 것은 ``무엇을 입력으로 받고, 무엇을 식별 대상으로 삼는가''다.
본 부록은 이론장(Chapter 1--5)의 결과를 \textbf{수정 없이} 입력으로 받아, 그 위에 \emph{피팅 방법론} 계층만
적층한다. 따라서 이 절은 (i) 이론장이 본 부록에 넘겨주는 식과 (ii) 본 부록이 데이터로 식별할 파라미터를 먼저
한 자리에 모은다 — 이 목록이 닫혀 있어야 ``자유 knob 의 무분별한 추가''라는 과적합의 첫 구멍이 막힌다.

아래 전달식 목록은 이론장 본문의 결과를 의미로 인용한 것이다(클린 이론장 확정 후 eq 라벨 재동기 대상).

\begin{linkbox}
\textbf{(L1) 전하 보존식(중심 상태식, 이론장 Ch1).} $Q_\cell\,q=Q_\bg(V_n,T)+\sum_{j=1}^{N_p}Q_{j,\tot}\xi_j$. $V_n$ 은
그 해(외부 입력 아님); 단조성 $\partial Q_\bg/\partial V_n\ge\epsilon_Q>0$.\\
\textbf{(L2) 진행률 동역학(이론장 Ch1$\cdot$Ch3).} $\dfrac{\dd\xi_j}{\dd t}=k_j[\xi_{\eq,j}(V_n,T)-\xi_j]$ 또는 forward/
backward 형 $r_j^+(1-\xi_j)-r_j^-\xi_j$; detailed balance $r_j^+/r_j^-=\xi_{\eq,j}/(1-\xi_{\eq,j})$;
$k_j=k_0(T)e^{-(G-\chi_j\mathcal A_j)/RT}$, $\mathcal A_j=s_{\phi,j}n_j^{\eff}F(V_{n,\drive}-U_j)$, 기본 $V_{n,\drive}=V_n$.\\
\textbf{(L3) 피팅 논리식과 평가 순서(이론장 Ch1).} $\dfrac{\dd Q}{\dd V_n}=\dfrac{C_\bg(V_n,T)}{1-Q_p\,(\dd\Theta_\mathrm{tail}/\dd Q)}$,
$C_\bg=\partial Q_\bg/\partial V_n$; 평가 순서 S1--S6 (charge-balance root $\to\xi_{\eq}\to k_j\to L(G)\to A_L\to$
kernel integral $\to\dd Q/\dd V_n\to$ 측정축 $\dd Q/\dd V_{n,\app}$).\\
\textbf{(L4) self-consistent 적분식(이론장 Ch1).} $\Theta(q)=\Theta_\mathrm{init}(q)+\int_0^q\mathcal K(q,q')\,
\Theta_\eq(V_n(q',\Theta(q')),T)\,\dd q'$, $\mathcal K(q,q')=\int_0^\infty A_L(L)\tfrac1L e^{-(q-q')/L}\dd L$ (Volterra 2종).\\
\textbf{(L5) closure 2-tier(이론장 Ch1).} \emph{직접 수치적분}(분포 $g$-grid)$=$항상 보존되는 기준 해법; \emph{Plan A}
$=$Fredholm ratio-substitution closed-form(사용자 PhD, \cite{lee2011,son2013,jcp2017})$=$검증 통과 시 가속 후보.\\
\textbf{(L6) 발열(이론장 Ch2$\cdot$Ch4).} $\dot{\mathcal Q}_\irr=\sum_j n_j^{\eff}I_j\eta_j\ge0$, $\dot{\mathcal Q}_\rev=
s^\conv|I|T\,\partial U/\partial T$ 류; calorimetry 부재 시 \textbf{forward prediction} 으로만(피팅 금지).\\
\textbf{(L7) 히스테리시스(이론장 Ch5).} signed current $I_s$$\cdot$$q_\stt$, 충전 부호 $s_{\phi,j}^b=-1$, branch target
$\xi_{\tar,j}^b=\xi_{\sstat,j}^b$, $\Delta V_\obs\ne\Delta V_\hys$, loop-area 분리 발열.\\
\textbf{(L8) 반증(이론장 Ch1).} N1 Arrhenius 직선성, N2 전위 의존 spectrum shift($\partial\ln L/\partial V_{n,\drive}<0$),
N3 저온 lineshape(평형 broadening vs disorder), N4 $\rho_G$ forward-only(역산 금지).
\end{linkbox}
\noindent\textbf{$\conv,\stt,\hys,\OCV$ 등 약어는 이론장 Ch1--5 정의 그대로}이며 본 부록은 이를 재정의하지 않는다.
이론장 본문 기본식은 수정하지 않는다. 본 부록은 이들을 \emph{데이터로 제약}하는 절차로 배치한다.

\textbf{피팅 파라미터의 전수 목록(이론장이 남긴 것).} 본 부록이 식별 대상으로 삼는 파라미터는 다음으로 닫힌다(이 외의
자유 knob 추가 금지). 각각 \emph{어느 장이 정의했고 어느 독립 실험이 1차로 제약하는지}를 \S\ref{sec:ch6_identif}
에서 매핑한다.
\begin{longtable}{p{0.20\textwidth} p{0.30\textwidth} p{0.44\textwidth}}
\toprule 파라미터 & 정의 장$\cdot$식 & 물리적 의미(피팅 시 해석) \\ \midrule \endhead
$U_j,\ w_j$ & Ch1 logistic & 전이 중심 전위$\cdot$평형 전이 폭(열역학 무대) \\
$Q_{j,\tot},\ Q_\bg(V_n,T)$ & Ch1 charge\_balance & 전이 용량 기여$\cdot$배경 누적 용량(전하 보존) \\
$\Delta H_{a,j},\ \Delta S_{a,j},\ k_0$ & Ch1 kact (\AL{10}) & 활성화 엔탈피$\cdot$엔트로피$\cdot$prefactor(Eyring) \\
$\chi_j(\equiv\beta_j$; 조건부, 대칭 Marcus$\cdot$정상 target 극한 한정$)$ & Ch1/Ch3 (\AL{11},31) & 배리어 낮춤 민감도 $=$ transfer coefficient \\
$\rho_G(g;T)$ & Ch1 spectrum (\AL{13}) & 배리어 확률밀도(\textbf{독립 입력}; 역산 금지, \AL{16}) \\
$R_n$ & Ch1 Vapp / Ch3 & 음극 apparent 분극 계수($\ne R_\ct\ne R_{n,\eff}$) \\
$k_{j,\lim}$ & Ch1 kphys & 수송$\cdot$확산$\cdot$유한자리 병목(독립 제약 시만) \\
$n_j^{\eff}=RT/(Fw_j)$ & Ch3 (\AL{34}) & effective 전자수(detailed balance 가 강제; 독립 자유 X) \\
$s_{\phi,j}^b,\ \xi_{\tar,j}^b,\ h_j$ & Ch5 (\AL{53}--56) & branch 부호$\cdot$target$\cdot$hysteresis state(충전 시) \\
\bottomrule
\end{longtable}

\textbf{의미.} 이 목록은 피팅 전 과정의 ``계정 과목표''다 — 이후 모든 절은 이 표의 어느 항을 \emph{어느 데이터로}
닫는지를 추적한다. 표에 없는 양을 자유화하려는 충동이 곧 과적합의 신호이며, 본 부록은 그 충동을 식별성
규칙(\S\ref{sec:ch6_identif})으로 차단한다. \emph{다음 절}은 이 목록 중 보유 데이터로 실제 제약 가능한 항과
곤란한 항을 먼저 가른다.

% ====================================================================
\section{보유 데이터와 식별 가능 / 불가 항목의 분리}\label{sec:ch6_data}
% ====================================================================
\textbf{질문.} 회귀를 돌리기 전에 물어야 할 것은 ``이 데이터로 \emph{무엇을} 제약할 수 있는가''다. 같은 곡선이라도
어떤 파라미터는 직접 결정되고 어떤 파라미터는 데이터가 침묵한다 — 후자를 자유화하면 곧장 과결정(over-fitting)으로
간다. 그래서 \textbf{피팅의 첫 단계는 회귀가 아니라 \emph{제약 가능성의 분리}}다.

1차 검증 기준은 사용자 보유 \textbf{음극 반쪽셀 GITT, STO--OCV 테이블, $0.1/0.2/0.5/1.0\,\mathrm{C}$ rate series}
다(\AL{63}). 이 데이터는 모델을 \emph{전부} 결정하지 않으므로, 제약 가능/불가 항을 먼저 나눈다.
\begin{longtable}{p{0.28\textwidth} p{0.34\textwidth} p{0.32\textwidth}}
\toprule 데이터 & 우선 제약 가능 & 직접 확정 곤란 \\ \midrule \endhead
반쪽셀 STO--OCV & $V_{n,\OCV}(q),U_j,w_j,Q_{j,\tot},Q_\bg$ 열역학 골격 & $k_j,\rho_G(g)$, transport limitation \\
반쪽셀 GITT & instantaneous polarization($R_n$ 일부), relaxation signature, OCV 정합 & $k_j$ 단독값, solid diffusion vs phase relaxation 완전 분리 \\
$0.1/0.2\,\mathrm C$ & 저율 quasi-equilibrium proxy, peak 위치/폭 & strict thermodynamic equilibrium \\
$0.5/1.0\,\mathrm C$ & kinetic lag, broadening, transport stress test & pure frozen kinetic limit \\
다온도 series & Arrhenius 기울기($E_a$), $\Delta S_a$ & 단일 온도만으로는 $T$-의존 분리 불가 \\
온도별 수명/full-cell & 정성적 열화 trend & $U_j(T),w_j(T),k_j(T),\Delta S_j$ 정량 분리 \\
\bottomrule
\end{longtable}

\textbf{의미.} 왼쪽 열은 ``이번 데이터가 직접 말해 주는 것'', 오른쪽 열은 ``추가 실험 없이는 침묵하는 것''이다 —
오른쪽 열을 무리하게 자유화하지 않는 것이 식별성 보존의 출발점이다. 데이터 자체가 \emph{어떤 제약을 줄 자격이
있는지}는 다음 두 box 가 정량 경계로 못박는다.

\begin{boundbox}
\textbf{데이터 제약(\AL{63}).} 단일 온도 자료로 $U_j(T),w_j(T),k_j(T)$ 의 온도 의존을 분리하지 않는다 — Arrhenius
기울기는 \emph{여러} 온도가 있어야 정의된다(\S\ref{sec:ch6_arrhenius}). calorimetry 없이 Ch2/Ch4 발열 파라미터를
피팅하지 않는다 — 발열 항은 \textbf{forward prediction}(미시$=$거시 $n^{\eff}I_j\eta_j$ 로 예측만)으로 두고
calorimetry 확보 후 검증한다(\AL{67}).
\end{boundbox}
\begin{boundbox}
\textbf{등온 가정 유효 범위.} 등온 피팅은 온도 상승이 충분히 작거나 외부 온도 제어가 강할 때만 유효하다. 고율
자기발열로 $T(t)$ 가 유의하게 변하면 $U_j(T),w_j(T),k_j(T),R_n(T),k_{j,\lim}(T),\rho_G(g;T)$ 가 함께 변하므로
등온 C-rate 자료만으로 kinetic parameter 를 확정하지 않고 Ch2/Ch4 thermal coupling 으로 돌아간다(\AL{67}).
\end{boundbox}

제약 가능 항이 정해졌으니, \emph{다음 절}부터는 이론장의 평가 순서 S1--S6 을 ``파라미터의 함수로 모델 ICA 를
한 번 계산하는'' 구체적 수치 연산으로 풀어, 그 계산을 데이터에 거는 피팅의 inner loop 를 만든다.

% ====================================================================
\section{피팅 평가 순서 S1--S6 의 구현 (이론장 평가 순서의 피팅 워크플로화)}\label{sec:ch6_S}
% ====================================================================
\textbf{질문.} 이론장의 결론식은 평가 순서 S1--S6 으로 \emph{한 점}의 $\dd Q/\dd V_n$ 을 계산한다. 그런데 피팅이
하려는 일은 ``이 계산을 파라미터의 함수로 두고, 측정 ICA 와 비교해 residual 을 최소화''하는 것이다. 그렇다면
S1--S6 의 각 단계는 \emph{어떤 수치 연산}으로 실현되는가? 이 수치 연산이 곧 피팅의 \textbf{inner loop}(주어진
파라미터에서 모델 ICA 를 한 번 평가하는 forward 계산)이다 — 본 절은 그 inner loop 를 단계별로 명시한다.

\subsection{S1 — 전하 보존식 root 로 $V_n$ 결정 (피팅 inner-loop 의 첫 연산)}\label{sec:ch6_S1}
\textbf{무엇을 구하나.} inner loop 의 첫 미지수는 내부 전위 $V_n$ 이다. 이론장이 강조하듯 $V_n$ 은 외부에서
읽는 값이 아니라 전하 보존식의 \emph{해}다. 주어진 진행 좌표 $q$ 와 진행률 $\bm\xi=(\xi_1,\dots,\xi_{N_p})$ 에서
$V_n$ 은 전하 보존식을 $V_n$ 에 대해 푸는 \emph{root finding} 으로 얻는다:
\begin{equation}
F_\cb^\dyn(V_n;q,\bm\xi)\;=\;Q_\bg(V_n,T)+\sum_{j=1}^{N_p}Q_{j,\tot}\xi_j-Q_\cell q\;=\;0.
\label{eq:ch6_dyn_root}
\end{equation}
$\bm\xi$ 는 이 root finding 에서 \emph{독립 입력}(시간 적분이 따로 갱신)이므로 Jacobian 은 배경 용량 기울기뿐이다:
\begin{equation}
\frac{\partial F_\cb^\dyn}{\partial V_n}=\frac{\partial Q_\bg}{\partial V_n}\;\ge\;\epsilon_Q>0.
\label{eq:ch6_dyn_deriv}
\end{equation}
이론장의 단조성 floor 가 유지되면 해 존재 범위에서 root 는 유일하다(중간값 정리 $+$ 단조성).

\textbf{이걸로 피팅에서 무엇을 하나(평형 OCV root 와의 구분).} 같은 root finding 이라도 평형/준평형에서는 모드가
다르다. 평형에서는 $\xi_j=\xi_{\eq,j}(V_n)$ 이므로 root 식이 달라지고 Jacobian 에 $\xi_{\eq}$ 항이 추가된다:
\begin{equation}
F_\cb^\OCV(V_n;q)=Q_\bg(V_n)+\sum_j Q_{j,\tot}\xi_{\eq,j}(V_n)-Q_\cell q=0,\qquad
\frac{\partial F_\cb^\OCV}{\partial V_n}=\frac{\partial Q_\bg}{\partial V_n}+\sum_j Q_{j,\tot}\frac{\partial\xi_{\eq,j}}{\partial V_n}.
\label{eq:ch6_ocv_root}
\end{equation}
두 root 의 derivative 가 다르므로(식~\eqref{eq:ch6_dyn_deriv} vs \eqref{eq:ch6_ocv_root}) 같은 Newton
Jacobian 으로 처리하면 안 된다 — 이것이 S1 을 피팅 inner loop 에 쓸 때의 핵심 주의다. 나아가
식~\eqref{eq:ch6_ocv_root} 의 $\partial F_\cb^\OCV/\partial V_n$ 은 이론장 Ch.2 의 평형 chemical
capacitance(OCV 온도계수 유도에 쓴 $\dd Q/\dd V_{n,\OCV}$)와 같은 양이므로, \textbf{S1 의 평형 극한이 Ch.2 의
평형 미분용량과 자동 정합}하는지가 곧 내부 정합 점검이 된다 — 둘이 어긋나면 피팅이 잘못된 것이다.
\begin{verifybox}
\textbf{root-find 종료 기준은 임의 상수가 아니다(\AL{64}, GS-4$'$).} Newton 종료는 ``residual 이 충분히 작다''가
아니라 \emph{전하 단위의 물리 scale} 로 정규화한다: $|F_\cb^\dyn|\le\varepsilon_Q^{\mathrm{root}}\equiv
\delta_q\,Q_\cell$, 여기서 $\delta_q$ 는 진행 좌표 $q$ 의 측정/계산 분해능(예: $\dd q$ 격자 간격)이다. 즉 ``배경
용량 잔차가 한 격자 스텝의 외부 전하보다 작다''는 \emph{측정 분해능 기준}이며, 무근거 숫자가 아니다. 부동소수점
한계상 $\varepsilon_Q^{\mathrm{root}}$ 을 기계 epsilon 아래로 두지 않는다($\sqrt{\varepsilon_{\mathrm{mach}}}$
근방이 Newton 의 통상 floor, \cite{brenan1996}).
\end{verifybox}

\subsection{S2 — 평형 target $\xi_{\eq,j}$ (logistic 평가)}\label{sec:ch6_S2}
\textbf{무엇을 구하나.} S1 의 $V_n$ 이 정해지면, 그 전위에서 각 전이의 \emph{평형 목표} 진행률 $\xi_{\eq,j}$ 를
얻는다. 이론장의 logistic 형태 $\xi_{\eq,j}(V_n,T)=[1+\exp(-s_{\phi,j}(V_n-U_j)/w_j)]^{-1}$ 를 그 전위에서 평가한다.
\textbf{이걸로 피팅에서 무엇을 하나.} 이 $\xi_{\eq,j}$ 가 다음 단계 동역학의 \emph{끌개}(target)이며, 그 모양을
정하는 $U_j,w_j$ 는 S1 의 $Q_{j,\tot},Q_\bg$ 와 함께 \emph{저속 OCV} 가 1차로 제약한다(\S\ref{sec:ch6_ocvfit}).
$n_j^{\eff}=RT/(Fw_j)$ 는 $w_j$ 가 정해지면 \emph{종속}이며 독립 자유 파라미터가 아니다(이론장 Ch3 \AL{34}).

\subsection{S3 — 속도상수 $k_j$ (Eyring 평가, 병목 직렬 합성)}\label{sec:ch6_S3}
\textbf{무엇을 구하나.} target 이 정해졌으니 ``그 target 으로 얼마나 빨리 다가가는가''를 정하는 속도상수 $k_j$ 가
필요하다. 구동력 $\mathcal A_j=s_{\phi,j}n_j^{\eff}F(V_{n,\drive}-U_j)$(기본 $V_{n,\drive}=V_n$)를 먼저 만들고,
$k_j^{\mathrm{act}}=k_0(T)e^{-(G-\chi_j\mathcal A_j)/RT}$ 를 평가한다(이론장 Eyring 형). 병목이 유효하면 이론장의
\emph{직렬 시간척도 합} $1/k_j^{\eff}=1/k_j^{\mathrm{act}}+1/k_{j,\lim}$ 로 합성한다.
\textbf{이걸로 피팅에서 무엇을 하나.} $k_j^{\eff}$ 는 S5 의 꼬리 길이와 inner-loop 시간 적분(\S\ref{sec:ch6_dae})에
직접 들어간다. \textbf{강구동에서 $k_j^{\mathrm{act}}$ 가 커지는 것을 임의 cap 으로 자르지 않는다}(GS-4); 수치
stiffness 는 \S\ref{sec:ch6_dae} 의 implicit solver 가 흡수한다.

\subsection{S4 — 배리어$\to$길이 $L(G)$ 와 spectrum $A_L$}\label{sec:ch6_S4}
\textbf{무엇을 구하나.} $k_j$ 가 정해지면 ``그 속도가 만드는 꼬리가 얼마나 긴가''를 길이 척도로 환산한다. 이론장의
닫힌 지수 매핑 $L(G)=\dfrac{|I|}{Q_\cell k_0(T)}\exp[(G-\chi_j\mathcal A_j)/RT]$ 로 배리어 $G$ 를 완화 길이 $L$ 로
보내고, 변수변환 밀도 보존 $A_L^{\mathrm{prob}}(L)=\rho_G(G(L))\,(RT/L)\,H(L-L_{\min})$ 로 배리어 분포를 길이
분포로 보낸다.
\textbf{이걸로 피팅에서 무엇을 하나.} $A_L(L)$ 이 다음 단계 kernel integral 의 가중치이며, 이것이 곧 ``여러 mode
가 겹쳐 만드는 관측 꼬리''의 재료다. \textbf{Heaviside $H$ 는 이론장 변수변환의 자코비안 support
하한($G\ge0\Leftrightarrow L\ge L_{\min}$)에서 \emph{유도}된 정의역 경계}이며, 피팅 편의를 위한 임의 절단이
아니다(GS-4 비위반). 수치적으로는 $g$-grid 의 하한을 $L_{\min}$ 에 맞추어 음의 배리어 영역을 적분에서 제외한다.

\subsection{S5 — 꼬리 중첩 $\dd\Theta_\mathrm{tail}/\dd q$ (kernel integral; 자기참조 풀이는 \S\ref{sec:ch6_grid})}\label{sec:ch6_S5}
\textbf{무엇을 구하나.} 분포 $A_L$ 위로 ``각 길이 mode 의 꼬리를 모두 더한'' 관측 꼬리의 증가율을 얻는다. 이론장의
kernel integral $\dfrac{\dd\Theta_\mathrm{tail}}{\dd q}=\int_0^\infty A_L(L)\,\tfrac1L\,e^{-(q-q_a)/L}\,\dd L$ 를
평가한다. \textbf{이걸로 피팅에서 무엇을 하나.} 이 양이 다음 단계 ICA 식의 분모를 채운다. 다만 $\xi_j$ 와 $V_n$ 이
전하 보존으로 얽혀 있어 이는 \emph{독립} 적분이 아니라 self-consistent Volterra 2종(L4)이며, 그 \emph{직접 수치
풀이가 본 부록 closure 의 기준 해법}이다(\S\ref{sec:ch6_grid}--\ref{sec:ch6_closure}에서 그 풀이를 구성한다).

\subsection{S6 — $\dd Q/\dd V_n$ 과 측정축 $\dd Q/\dd V_{n,\app}$}\label{sec:ch6_S6}
\textbf{무엇을 구하나.} 마지막으로 한 점의 ICA 값을 얻는다. 이론장 결론식 $\dd Q/\dd V_n=C_\bg/(1-Q_p\,\dd\Theta/\dd Q)$
를 평가하고, 측정축으로는 apparent 전위 $V_{n,\app}=V_n+s_I|I|R_n$ 로 환산한다: 등온 정전류에서
$\dd V_{n,\app}/\dd q=\dd V_n/\dd q+s_I|I|\,\partial R_n/\partial q$, DVA 는 그 역수다.
\textbf{이걸로 피팅에서 무엇을 하나.} \textbf{측정 ICA 는 $\dd Q/\dd V_{n,\app}$ 와 비교}하며, 모델 출력의
$V_n$-축을 직접 데이터에 맞추지 않는다 — $R_n$ 이 축을 이동시키기 때문이고, 이 축 이동이 곧 식별성 문제의
씨앗이다(\S\ref{sec:ch6_identif}). 이로써 S1$\to$S6 의 inner loop 한 바퀴가 한 점의 모델 ICA 를 내놓는다.

\begin{practicebox}
\textbf{inner-loop 안정화(본문 물리식 아님, GS-4$''$).} 초기값 탐색 단계에서 $w_j>0$, $Q_{j,\tot}>0$,
$0\le\chi_j\le1$ 같은 물리 제약을 만족시키려 $\log$/$\mathrm{logit}$ 같은 bounded reparametrization 을 쓸 수
있다. 이는 \emph{최적화 변수의 변환일 뿐} 모델 물리식이 아니며, 최종 보고에서는 원 파라미터값으로 되돌려
신뢰구간과 함께 제시한다. rate 폭주를 임시 clip 하는 것도 디버깅 한정이며 최종 해석에서 제거한다(물리 병목
$k_{j,\lim}$ 으로 대체).
\end{practicebox}

S1--S6 이 한 점을 내놓았으니, \emph{다음 절}은 곡선 \emph{전체}를 얻기 위해 $\xi_j(t)$ 의 시간 전개를 푸는 수치
엔진을 구성한다.

% ====================================================================
\section{피팅 inner-loop 의 수치 엔진 (root-constrained ODE / index-1 DAE)}\label{sec:ch6_dae}
% ====================================================================
\textbf{질문.} 피팅이 한 파라미터 집합에서 모델 ICA 곡선 \emph{전체}를 얻으려면 한 점이 아니라 $\xi_j(t)$ 의 시간
전개를 풀어야 한다. 그렇다면 이 시간 전개는 어떤 종류의 방정식인가? 정전류 음극 방전에서 $\dd q/\dd t=|I|/Q_\cell$
이고, 상태변수는 $\xi_j$, $V_n$ 은 전하 보존식이 매 시점 정하는 \emph{algebraic variable} 다. 즉 미분식 1조
$+$ 대수 제약 1조의 \textbf{index-1 DAE}(\AL{60}, \cite{brenan1996,hindmarsh2005})다:
\begin{align}
\frac{\dd\xi_j}{\dd t}&=k_j^{\eff}(V_n,q,|I|,T)\,[\xi_{\eq,j}(V_n)-\xi_j],\label{eq:ch6_ode}\\
0&=F_\cb^\dyn(V_n;q,\bm\xi).\label{eq:ch6_alg}
\end{align}
\textbf{왜 index-1 인가(무비약).} DAE 의 index 는 대수 제약을 미분해 ODE 로 환원하는 데 필요한 미분 횟수다.
식~\eqref{eq:ch6_alg} 을 시간으로 한 번 미분하면 $\dfrac{\partial Q_\bg}{\partial V_n}\dot V_n+\sum_j Q_{j,\tot}
\dot\xi_j-Q_\cell\dot q=0$ 이고, $\partial Q_\bg/\partial V_n\ge\epsilon_Q>0$(식~\eqref{eq:ch6_dyn_deriv})이라
$\dot V_n$ 에 대해 \emph{한 번}에 풀린다 — 따라서 index 가 정확히 1 이다(\AL{60}).
\textbf{이걸로 피팅에서 무엇을 하나.} 단조성 floor 가 이 풀림의 열쇠이므로, 피팅 중 $\partial Q_\bg/\partial V_n$
이 0 에 접근하면 DAE 가 ill-conditioned 가 되고 이는 \emph{파라미터가 비물리}라는 신호다(reject; \S\ref{sec:ch6_S1})
— 즉 수치 조건수 자체가 파라미터 타당성의 진단기가 된다.

\textbf{시간 적분 절차(정전류; $q(t)$ 명시이므로 root-constrained ODE 로 구현).} 위 DAE 를 한 스텝 전진시키는
절차는 S1--S6 의 inner loop 를 시간 루프로 감싼 것이다:
\begin{enumerate}[label=T\arabic*),leftmargin=2.4em]
\item 현재 $q(t),\xi_j(t)$ 에서 식~\eqref{eq:ch6_dyn_root} 로 $V_n(t)$ root-find(S1).
\item $V_n(t)$ 에서 $\xi_{\eq,j}$(S2), $k_j^{\mathrm{act}},k_{j,\lim},k_j^{\eff}$ 또는 $r_j^\pm$(S3) 계산.
\item 식~\eqref{eq:ch6_ode}(또는 forward/backward 형)로 $\dd\xi_j/\dd t$ 계산.
\item implicit solver(BDF 또는 Radau)로 다음 시간으로 진행(\AL{60}, \cite{hindmarsh2005}).
\item 휴지 구간에서는 $q$ 고정, $\xi_j$ relaxation 과 전하 보존 root 를 함께 갱신(Ch2/Ch5 rest; 내부 $I_j\ne0$).
\end{enumerate}
\begin{boundbox}
\textbf{stiff system 처리(cap 아님, GS-4).} 강구동에서 $k_j^{\mathrm{act}}$ 또는 $r_j^+$ 가 매우 커질 수 있다.
이를 arbitrary cap 으로 자르지 않고 \emph{stiff system} 으로 취급한다(implicit BDF/Radau 가 큰 $k_j$ 를 안정적으로
처리; \cite{brenan1996,hindmarsh2005}). 실제 rate 제한은 $k_{j,\lim}$, transport, 유한 전류, site availability
같은 물리적 병목으로 명시한다(이론장 Ch3 정합). 즉 \textbf{수치 stiffness 의 해법은 solver 선택이지 모델
변형이 아니다}.
\end{boundbox}
\begin{verifybox}
\textbf{solver tolerance 의 근거(\AL{64}, GS-4$'$).} BDF/Radau 의 상대$\cdot$절대 tolerance $(\mathrm{rtol},
\mathrm{atol})$ 는 임의 상수가 아니라 상태변수의 물리 scale 로 정한다: $\xi_j\in[0,1]$ 이므로 $\mathrm{atol}$ 은
ICA 의 관심 분해능(예: peak 높이 대비 꼬리 기여의 유의 자릿수)으로, $\mathrm{rtol}$ 은 측정 잡음 대비 over-resolve
를 피하는 수준으로 둔다. solver tolerance 를 데이터 잡음보다 훨씬 작게 두는 것은 계산 낭비이고, 훨씬 크게 두면
꼬리(tail-dominated, 작은 기여)를 놓친다 — \textbf{tolerance 는 측정 잡음 floor 와 ICA 분해능 사이}에 둔다.
\end{verifybox}

시간 적분 절차가 갖춰졌으니, \emph{다음 절}은 그 안에 들어가는 핵심 부품 — self-consistent Volterra 의 직접 수치
풀이(기준 해법) — 를 구성한다.

% ====================================================================
\section{closure: self-consistent Volterra 의 직접 수치 풀이 (기준 해법)}\label{sec:ch6_grid}
% ====================================================================
\textbf{질문.} 이론장의 spectrum 은 배리어 분포 $\rho_G$ 가 만드는 \emph{연속} relaxation-length 분포다. 그러나
컴퓨터는 연속 분포를 그대로 적분하지 못한다 — 그렇다면 이 연속 분포를 \emph{어떻게} 유한 계산으로 푸는가? 가장
정직한 답은 분포를 유한 격자로 이산화하는 것이다. 각 배리어 격자점 $g_m$ 의 집단이 독립 1차 완화하고, 전하
보존이 매 시점 $V_n$ 으로 결합한다:
\begin{equation}
\frac{\dd\xi_j(g_m,t)}{\dd t}=k_j^{\eff}(g_m,V_n,q,|I|)\,[\xi_{\eq,j}(V_n)-\xi_j(g_m,t)],\qquad
\xi_j(t)=\sum_{m=1}^{N_g}w_m^{(g)}\,\rho_G(g_m)\,\xi_j(g_m,t),
\label{eq:ch6_grid}
\end{equation}
$\sum_m w_m^{(g)}\rho_G(g_m)\to1$ (격자 정규화; quadrature 가중 $w_m^{(g)}$).
\textbf{이걸로 피팅에서 무엇을 하나.} \textbf{이 직접 $g$-grid 적분이 본 부록의 기준 해법}이다 — 주어진
$\rho_G$ 와 완화 closure 하에서 \emph{이산화 오차만 남는 numerically exact reference} 이며, 이후 모든 축약
해법(Plan A 포함)은 \emph{이 해에 대해} 검증한다. 이는 이론장 self-consistent Volterra(L4)의 직접 수치 평가다.
\begin{verifybox}
\textbf{이중 계상 점검(이론장 boundbox 의 수치 재확인).} 상수 목표 극한($\Theta_\eq=\Theta_0$, 단일 mode
$A_L=\delta(L-L_0)$, $\Theta_\mathrm{init}=0$)에서 식~\eqref{eq:ch6_grid} 의 이산 해는 $\Theta(q)=\Theta_0
(1-e^{-q/L_0})$ 로 이론장 single-mode ODE 와 일치하고 $q\to\infty$ 서 $\Theta_0$ 로 수렴해야 한다. 만약 kernel 을
\emph{목표} 대신 \emph{잔차}에 곱하면 정상상태가 $\Theta_0/2$ 로 어긋난다 — 수치 구현이 이론장 Volterra(목표에
곱)와 같은지 이 극한으로 검증한다.
\end{verifybox}
\begin{boundbox}
\textbf{격자 분해능과 stretched tail(\AL{62}).} 이론장의 꼬리는 \emph{tail-dominated}(긴 $L$ mode 가 지배)다.
$g$-격자가 큰 $L$(작은 $k_j$, 큰 배리어) 영역을 충분히 덮지 않으면 꼬리를 놓친다. 따라서 격자 상한은 관심
온도/전위에서 $L_{\max}$ 가 관측 윈도를 덮도록 잡고, 격자 수 $N_g$ 는 $\Theta(q)$ 가 $N_g$ 에 대해 수렴할 때까지
늘린다(자기 수렴 점검; 임의 고정값 아님).
\end{boundbox}
\begin{flagbox}
\textbf{독립 완화 closure = 이론장이 본 부록에 위임한 FLAGGED 항목(미이행).} 식~\eqref{eq:ch6_grid} 의 ``각 $g_m$
집단이 \emph{독립} 1차 완화''는 이론장 평균장(mean-field) ansatz 다. 이론장은 이를 flagbox 로 정직 표기하며
``흑연 staging 은 nucleation$\cdot$상경계 이동을 동반하는 \emph{협동적} 1차 상전이라 mode 간 결합(Avrami)을 무시한
평균장이며, 그 타당성은 FLAGGED — 본 부록(\S\ref{sec:ch6_reduction})에서 검증''이라 위임했다. 그러나 위임된
\emph{협동적 staging 검증 자체는 본 작업에서 미이행}이다(독립 완화의 정량 타당성은 operando/calorimetry 측정 +
협동 모형 비교를 요하며 본 이론 범위 밖). 따라서 본 절 $g$-grid 의 ``numerically exact reference''는 그 closure
\emph{하에서만} exact 이며 \emph{exactness 는 이산화 오차에 한하고 closure 자체에는 미치지 않는다} — FLAGGED 로
승계한다.
\end{flagbox}

기준 해법이 확정됐으니, \emph{다음 절}은 ``기준 해 대비 검증을 전제로'' 허용되는 수학적 가속을 정리한다.

% ====================================================================
\section{closure 축약: 허용되는 수학적 가속과 그 검증}\label{sec:ch6_reduction}
% ====================================================================
\textbf{질문.} 직접 격자(\S\ref{sec:ch6_grid})는 정확하지만 비용이 커, 피팅 반복마다 매번 풀기엔 부담이다. 그렇다면
\emph{물리 의미를 훼손하지 않으면서} 가속하는 방법은 무엇이고, 어떤 조건에서 허용되는가? 핵심 규칙은 하나다 —
\emph{기준 해 대비 검증}돼야 한다. 다음 축약은 수학적 근거가 있을 때 허용한다.
\begin{enumerate}[label=\textbf{(\alph*)},leftmargin=2.4em]
\item \textbf{Adaptive quadrature.} $\rho_G(g)$ 가 매끄럽고 support 제한($L\ge L_{\min}$)/tail 빠른 감소면, 모델을
  바꾸지 않고 적분 정확도만 확보하는 수치 적분법이므로 허용(\AL{65}).
\item \textbf{Moment matching.} 분포 폭이 작고 kernel 이 매끄러우면 $k_j(g)$/memory kernel 을 저차 moment 로 근사.
  단 moment closure 가 \emph{tail-dominated} kinetics 를 놓치지 않는지 기준 해와 비교한다(이론장 stretched tail 주의).
\item \textbf{Prony/rational approximation.} 연속 relaxation spectrum 을 유한 exponential mode 로 근사(\AL{65}):
\begin{equation}
K_j(t)=\int_0^\infty\rho_G(g)\,k_j(g)\,e^{-k_j(g)t}\,\dd g\;\approx\;\sum_{\ell=1}^{N_r}a_\ell\,e^{-t/\tau_\ell},
\qquad a_\ell\ge0,\ \tau_\ell>0.
\label{eq:ch6_prony}
\end{equation}
\end{enumerate}
\textbf{$a_\ell\ge0,\tau_\ell>0$ 의 근거(무비약).} $K_j(t)$ 는 \emph{양의} 분포 $\rho_G k_j$ 위의 지수 감쇠 평균이라
완전 단조(completely monotone) 함수다; 그 유한 mode 근사가 물리적 완화이려면 각 mode 가 \emph{음이 아닌} 진폭과
\emph{양의} 시간상수를 가져야 한다(음수 가중/허수 time constant 는 진동/발산을 낳아 비물리). 따라서 부호 제약은
피팅 편의가 아니라 \emph{완전 단조성의 보존}이다(\AL{65}; KWW stretched 함수와 완화시간 분포의 등가성
\cite{lindsey1980}, 지수 감쇠의 연속합 \cite{johnston2006}).
\textbf{이걸로 피팅에서 무엇을 하나.} 정전류에서 시간$\leftrightarrow$좌표 대응은 $t-t'\leftrightarrow(q-q')Q_\cell/|I|$,
$\tau_\ell\leftrightarrow L_\ell$ 이므로 식~\eqref{eq:ch6_prony} 는 이론장 relaxation-length spectrum 의 유한 mode
표현이며, 검증을 통과하면 inner loop 의 kernel 평가를 닫힌 합으로 가속한다.
\begin{boundbox}
\textbf{축약의 한계(되먹임).} $V_n$ 이 전하 보존으로 $\xi_j$ 와 결합돼 고율/큰 lag 에서 $\xi_j\to V_n\to\xi_{\eq,j}
\to\dd\xi_j/\dd t$ 되먹임이 강해진다. Prony/rational/moment 축약은 \emph{준정상$\cdot$약결합} 근사로만 쓰고, 직접
$g$-grid DAE 대비 오차를 \S\ref{sec:ch6_closure} 의 절차로 측정한다. 되먹임 오차가 허용 범위를 넘으면 기준
해법으로 회귀한다.
\end{boundbox}

위 (c) Prony 와 같은 ``닫힌 형태 가속''의 대표가 사용자 PhD 의 Plan A 다. \emph{다음 절}은 그 Plan A 의 위치와,
그것을 채택할지 가르는 검증 tolerance 를 정직하게 정의한다.

% ====================================================================
\section{Plan A (Fredholm ratio-substitution closed-form) 의 위치와 ★ 검증 tolerance 의 정직한 정의}\label{sec:ch6_closure}
% ====================================================================
\textbf{질문.} 사용자 PhD 의 \textbf{Plan A}(ratio-substitution closed-form, \cite{lee2011,son2013,jcp2017})는
문헌적$\cdot$수학적 기반이 있는 \emph{가속 후보}다(\AL{61}). 그렇다면 이것을 기본 모델로 써도 되는가? 답은 ``아니오,
검증 대상으로 둔다''이다 — \textbf{직접 $g$-grid 기준 해를 대체하는 기본 모델이 아니라 검증 대상}이며, 이론장의
승격 조건 M1--M5 를 본 부록에서 수치적으로 집행한다.
\begin{enumerate}[label=M\arabic*),leftmargin=2.4em]
\item \textbf{kernel mapping.} 흑연 배리어 분포 문제로의 kernel mapping 이 명확하고 적분 안팎 function class 가 같다.
\item \textbf{function class.} ratio-substitution 이 가정하는 분포 class 가 흑연 $\rho_G$ 와 양립한다.
\item \textbf{수렴성.} ratio truncation/Neumann series 의 수렴이 수치적으로 확인된다.
\item \textbf{$\delta$ baseline.} $\rho_G(g)\to\delta(g-g_0)$ 극한에서 단일 배리어 해(이론장 single-mode
  $r\propto e^{-(q-q_a)/L_0}$)로 환원된다.
\item \textbf{정량 오차.} 기준 해 대비 상대오차가 검증 tolerance 이하다(아래 ★ 정의).
\end{enumerate}
M5 가 가리키는 ``검증 tolerance''는 본 부록에서 가장 조심해야 할 양이다 — \emph{다음 소절}이 이를 정직하게
재정의한다.

\subsection{★ 검증 tolerance $\varepsilon_\tol$ 의 정직한 재정의 (FLAGGED 허위 정정)}\label{sec:ch6_epsilon}
\textbf{문제(적발).} 구 consolidated 판의 M5 는 정량 오차 조건을
\[
\varepsilon=\frac{\lVert\Theta_{\mathrm A}-\Theta_{\mathrm B}\rVert}{\lVert\Theta_{\mathrm B}\rVert}\le\varepsilon_\tol
\]
로 쓰면서 $\varepsilon_\tol$ 을 \emph{established 임계값처럼} 사용했다. 그러나 (i) 그 \emph{값}의 근거가 본문에
없고, (ii) 문서 전체에 정직 표기용 flagbox 가 0건이었다. 이는 GS-3$\cdot$GS-4 위반(근거 없는 수치 임계값)이다.
본 소절은 이를 \textbf{근거 있는 수치 tolerance 로 재서술}하되, 정당화되지 않는 특정 숫자는 FLAGGED 로 정직 표기한다.

\textbf{무엇을 비교하는가(정의 명확화).} $\Theta_{\mathrm B}$ 는 직접 $g$-grid 기준 해(\S\ref{sec:ch6_grid}),
$\Theta_{\mathrm A}$ 는 Plan A closed-form 해다. norm $\lVert\cdot\rVert$ 은 \emph{관심 tail-window} 위의 이산
$L^2$ 노름으로 명시한다: $\lVert f\rVert^2=\sum_{i\in\mathcal W}(f(q_i))^2\,\Delta q$, $\mathcal W=$ post-peak 꼬리
구간. window 를 명시하지 않으면 peak 영역의 큰 값이 분모를 지배해 꼬리 오차를 가린다(\textbf{꼬리가 본 이론의
관심사}이므로 window 가 필수다).

\textbf{tolerance 의 \emph{근거 있는} 결정(임의 상수 금지).} $\varepsilon_\tol$ 은 ``작으면 좋다''가 아니라 두 물리
scale 의 \emph{큰 쪽}으로 둔다(\AL{66}):
\begin{equation}
\varepsilon_\tol\;=\;\max\!\Big(\underbrace{\varepsilon_{\mathrm{disc}}}_{\text{기준 해 자체의 이산화 오차}},\
\underbrace{\varepsilon_\noise}_{\text{측정 잡음으로 인한 ICA 불확실도}}\Big).
\label{eq:ch6_epsilon}
\end{equation}
\emph{한 줄씩 근거}: (a) Plan A 를 기준 해 $\Theta_{\mathrm B}$ 보다 더 정확하게 만들 수는 없다 — $\Theta_{\mathrm B}$
자신이 격자 이산화 오차 $\varepsilon_{\mathrm{disc}}$ 를 품기 때문이다. 따라서 $\varepsilon_\tol<\varepsilon_{\mathrm{disc}}$
를 요구하는 것은 무의미하다(기준 해의 정밀도를 넘는 일치를 강요). $\varepsilon_{\mathrm{disc}}$ 는 \S\ref{sec:ch6_grid}
의 격자 자기수렴(격자 배가 시 $\Theta_{\mathrm B}$ 변화)으로 \emph{측정}한다. (b) 모델 예측을 데이터 잡음보다 더
정밀하게 일치시키는 것은 관측 가능한 차이가 아니다 — ICA 의 잡음 유발 불확실도 $\varepsilon_\noise$ 는 측정
전압/전류 잡음을 $\dd Q/\dd V$ 로 전파해 추정한다($dQ/dV$ 는 미분이라 잡음이 증폭됨, \cite{dubarry2022,fly2020}
\AL{63}). 따라서 \textbf{Plan A 가 기준 해와 ``측정으로 구별 불가능한 만큼'' 일치하면 충분}하며, 그 경계가
식~\eqref{eq:ch6_epsilon} 다. 두 scale 중 \emph{큰 쪽}을 쓰는 이유는, 더 작은 쪽까지 맞추려는 것이 위 (a)(b)
어느 한쪽에서 무의미해지기 때문이다.
\begin{flagbox}
\textbf{$\varepsilon_\tol$ 의 \emph{특정 숫자값}은 미확정(\AL{66}).} 식~\eqref{eq:ch6_epsilon} 은 $\varepsilon_\tol$
을 \emph{측정 가능한 두 scale} 로 환원했으나, $\varepsilon_{\mathrm{disc}},\varepsilon_\noise$ 의 \emph{수치}는
사용자 보유 데이터의 실제 잡음 수준$\cdot$격자 수렴 거동에서 산출돼야 한다 — 본 이론 문서가 ``$\varepsilon_\tol=10^{-2}$''
같은 \emph{고정 숫자}를 established 로 제시하지 않는다. 구 consolidated 판이 $\varepsilon_\tol$ 을 무근거 임계값처럼
쓴 것을 본 소절이 \emph{측정 기반 정의}로 교체했고, 구체 값은 데이터 확보 후 결정하는 \textbf{미확정 항}으로 정직
표기한다(GS-3$\cdot$GS-4$'$ 준수). $L^2$ vs $L^\infty$ norm 선택, window 경계 정의도 데이터 의존이라 함께 미확정이다.
\end{flagbox}
\begin{enumerate}[label=F\arabic*),leftmargin=2.4em,start=6]
\item 실패 시 fallback 순서: 직접 grid(\S\ref{sec:ch6_grid}) $\to$ adaptive quadrature $\to$ Prony/rational
  (\S\ref{sec:ch6_reduction}). 즉 \emph{기준 해는 언제나 직접 grid} 이며 Plan A 는 통과 시에만 가속용으로 끼운다.
\end{enumerate}
\textbf{의미.} 이렇게 정의한 $\varepsilon_\tol$ 은 ``Plan A 를 가속용으로 끼워도 되는가''의 합격선이자, 끼우지
못할 때 직접 grid 로 회귀하는 트리거다. 다만 Plan A 가 \emph{가장 부정확해지는} 영역이 하필 본 이론의 관심
영역이라는 점을 box 가 경고한다.
\begin{boundbox}
\textbf{Plan A 의 broad-kernel 한계(★ \AL{62}, jcp2017 자기명시).} 사용자 논문(\cite{jcp2017}, Eq.~34 직후)은
``ratio-substitution 근사의 정확도는 reaction zone(분포)이 매우 넓어질 때 나빠진다''고 \emph{스스로 경고}한다.
graphite 의 넓은 spectrum($=$저온 stretched tail, \emph{관심 영역})은 바로 그 broad-kernel 영역이라 Plan A 가
\textbf{가장 부정확할 수 있다}. 또한 사용자 논문은 \emph{Fredholm}(고정 구간) ratio closure 이고 graphite 문제는
\emph{Volterra}(인과 상한 $q$, 이론장 L4)라 자기참조 선형화가 선행돼야 한다. 따라서 \textbf{Plan A 단독 채택 금지;
항상 직접 grid 대비 식~\eqref{eq:ch6_epsilon} 검증.} graphite transition 개수$\cdot$분포 폭$\cdot$hysteresis branch
offset 을 Fredholm correction factor 안에 흡수하지 않는다(식별성 보존).
\end{boundbox}

closure(수치 풀이)가 닫혔으니, \emph{다음 절}부터는 피팅의 또 다른 축 — 파라미터 식별성 — 으로 옮겨 간다.
같은 곡선을 만드는 파라미터들을 어떻게 분리하는가가 관건이다.

% ====================================================================
\section{파라미터 식별성(identifiability)과 공선형 차단}\label{sec:ch6_identif}
% ====================================================================
\textbf{질문.} closure 가 정확해도, \emph{여러 파라미터가 같은 곡선 변화를 만들어 서로 흉내 내면} 피팅은 유일한
답을 못 준다 — 이 \emph{공선형(co-linearity)}이 피팅의 핵심 위험이다. 그렇다면 흑연 ICA 에서 가장 위험한 공선형
쌍은 무엇이고, 그것을 어떻게 닫는가? 본 절은 가장 위험한 두 쌍을 먼저 닫고 일반 규칙을 제시한다.

\subsection{$R_n$ vs $\chi_j$ — 분극과 mobility 가속의 공선형}\label{sec:ch6_RnChi}
\textbf{왜 위험한가(둘 다 꼬리를 늘린다).} 분극 $R_n$ 은 apparent 전위축을 이동시켜(이론장 Vapp) ICA peak 를
넓히고, mobility 가속 $\chi_j$ 는 $L=|I|/(Q_\cell k_j)$ 를 통해 꼬리 길이를 직접 바꾼다(이론장 ideal-width 형
$\partial\ln L/\partial V_{n,\drive}=-\chi_j s_{\phi,j}F/RT$; 일반형은 $-\chi_j s_{\phi,j}n_j^{\eff}F/RT
=-\chi_j s_{\phi,j}/w_j$, $n_j^{\eff}=RT/Fw_j$). 한 데이터에서 둘을 동시에 자유 피팅하면 공선형이라 분리 불가다.

\textbf{분리 전략(순차 제약) — 이걸로 피팅에서 무엇을 하나.} \emph{독립 실험으로 $R_n$ 을 먼저 고정}한다:
펄스/전류 차단(GITT 의 순간 step)은 $|I|\to0$ 직후 전압 도약으로 \emph{순간} 분극을 주며, 이는 mobility
완화(느린 spectrum)와 시간척도가 분리된다. $R_n$ 을 이렇게 제한한 \emph{뒤에} 같은 모델로 $\chi_j$ 를 본다.
순서를 바꾸면(먼저 $\chi_j$ 자유화) 분극이 $\chi_j$ 로 흡수돼 가짜 가속을 얻는다.
\begin{keybox}
\textbf{세 저항은 다른 물리다(이론장 Ch1$\cdot$Ch3$\cdot$Ch4 계승).} apparent voltage-axis 계수 $R_n$ $\ne$ charge-transfer
저항 $R_\ct=RT/(FA_\eff j_{0,n})$(Ch3 small-signal) $\ne$ heat-generation 저항 $R_{n,\eff}$(Ch4 발열). 피팅에서
이 셋을 한 ``$R$''로 합치면 식별성과 발열 예측이 모두 무너진다 — \textbf{$R_n\ne R_\ct\ne R_{n,\eff}$} 를
파라미터 단계에서 분리 유지한다.
\end{keybox}

\subsection{$R_n/R_\ct/R_{n,\eff}$ 의 실험적 분리 (이론장 Ch3 계승)}\label{sec:ch6_threeR}
\textbf{무엇을 구하나.} keybox 가 구분한 세 저항을 \emph{각각 어느 실험이} 독립 제약하는지 못박는다. $R_\ct$ 는
\emph{소진폭 EIS/Tafel}(small-signal 선형 영역, \cite{bardfaulkner})에서, $R_n$ 은 \emph{ICA 전압축
이동}(apparent), $R_{n,\eff}$ 는 \emph{calorimetry}(발열률 $\div I^2$)에서 각각 독립 제약한다.
\textbf{이걸로 피팅에서 무엇을 하나.} 한 실험이 셋을 모두 주지 않으므로, 같은 데이터에서 동시에 자유화하지
않는다(\AL{39},\AL{47},\AL{68}) — 이 ``실험별 귀속''이 다음 소절 일반 규칙의 한 사례다.

\subsection{일반 식별성 규칙 — 동시 자유화 금지 집합}\label{sec:ch6_noFree}
\textbf{무엇을 구하나.} 위 두 사례를 일반화해, 같은 자료에서 동시에 자유화하면 식별성이 무너지는 파라미터 집합을
하나로 모은다(이론장 Ch1--3 누적 경고의 통합):
\begin{equation}
\{\,R_n,\ k_j^{\eff},\ w_j,\ \rho_G(g),\ Q_\bg,\ Q_{j,\tot},\ k_{j,\lim}\,\}.
\label{eq:ch6_no_free}
\end{equation}
\textbf{순차 제약 사슬(deliverable) — 이걸로 피팅에서 무엇을 하나.} 식~\eqref{eq:ch6_no_free} 의 공선형은
\emph{독립 실험을 순서대로 거는} 것으로만 풀린다. 본 부록의 권고 순서는 \S\ref{sec:ch6_ocvfit}--\ref{sec:ch6_crate}
에서 단계별로 실현한다:
\[
\underbrace{\text{저속 OCV}}_{U_j,w_j,Q_{j,\tot},Q_\bg}\ \to\ \underbrace{\text{GITT}}_{R_n,\ \text{완화 signature}}\ \to\
\underbrace{\text{다온도 Arrhenius}}_{\Delta H_a,\Delta S_a,k_0}\ \to\ \underbrace{\text{C-rate series}}_{\chi_j,\ k_{j,\lim}}.
\]
$\rho_G$ 는 어느 단계에서도 관측 꼬리에서 \emph{역산하지 않고}(비유일, \AL{16}) 독립 입력으로만 둔다.

\textbf{다리.} 이 사슬의 각 고리를 \emph{다음 네 절}이 한 단계씩 데이터로 실현한다 — 저속 OCV(\S\ref{sec:ch6_ocvfit})
$\to$ GITT(\S\ref{sec:ch6_gittfit}) $\to$ 다온도 Arrhenius(\S\ref{sec:ch6_arrhenius}) $\to$ C-rate(\S\ref{sec:ch6_crate}).

% ====================================================================
\section{S1 실현: 저속 OCV 로 열역학 골격 제약}\label{sec:ch6_ocvfit}
% ====================================================================
\textbf{질문(사슬 1번째 고리).} 먼저 무엇을 고정하는가? 가장 ``평형에 가까운'' 데이터로 열역학 골격
$U_j,w_j,Q_{j,\tot},Q_\bg$ 를 닫는다. 저속(준평형) OCV/STO--OCV 테이블을 평형 음극 전위 $V_{n,\OCV}(q)$ 의 실측
proxy 로 쓴다. \textbf{이 자료만으로 $U_j,w_j,Q_{j,\tot},Q_\bg$ 가 유일 결정되지는 않으므로} 다음 물리 제약 하에서
피팅한다(과적합 차단):
\begin{enumerate}[label=O\arabic*),leftmargin=2.4em]
\item 총용량 정합: $\sum_j Q_{j,\tot}+\Delta Q_\bg=Q_\cell\cdot\Delta q_{\mathrm{full}}$(이론장 전하 보존의 적분형).
\item $Q_\bg(V_n)$ 단조성 $\partial Q_\bg/\partial V_n\ge\epsilon_Q>0$ 및 smoothness(이론장 monotone).
\item transition 개수 $N_p$ 를 필요 이상 늘리지 않음(parsimony; 관측 $N_p\approx3\sim4$, \AL{3}).
\item $Q_{j,\tot}>0$(방전 방향 용량 기여; 이론장 Ch1).
\item $w_j$ 는 \emph{평형} transition width 이며 kinetic 배리어 분포 $\rho_G(g)$ 와 혼동 금지(이론장 구분; 단
  $n_j^{\eff}w_j=RT/F$ 정합, Ch3 \AL{34}).
\item 잔차를 임의 background wiggle 로 모두 흡수하지 않음(GS-4; 잔차는 다음 단계가 설명할 신호일 수 있음).
\end{enumerate}
\textbf{평형 형태의 판별(이론장 N1 의 실현) — 이걸로 무엇을 하나.} 저속 OCV 의 near-peak 감쇠를
$\log(\dd Q/\dd V)$ 대 $(V-U)$(직선이면 logistic, 이론장 logistic) 또는 대 $(V-U)^2$(직선이면 Gaussian-cumulative)로
그려 평형 isotherm 형태를 \emph{판별}한다. \textbf{logistic 이외 형태(erf 등)는 흑연에 대한 열역학적 유도가 없는
FLAGGED ansatz}(이론장 flagbox)이므로, 데이터가 logistic 을 기각하면 그때 BOUNDED ansatz 로만 도입한다.
\begin{boundbox}
\textbf{저속이 평형은 아니다.} $0.1/0.2\,\mathrm C$ 는 \emph{quasi-equilibrium proxy} 이지 strict thermodynamic
equilibrium 이 아니다(\AL{63}). peak 위치/폭에 잔존 kinetic broadening 이 섞일 수 있으므로, OCV 골격은 GITT
relaxation 끝점(\S\ref{sec:ch6_gittfit})과 교차 검증한다.
\end{boundbox}

골격이 고정됐으니, \emph{다음 절}은 사슬 2번째 고리 — GITT 로 분극과 완화 signature 를 제약 — 로 넘어간다.

% ====================================================================
\section{S1--S2 보강: GITT 로 분극과 완화 signature 제약}\label{sec:ch6_gittfit}
% ====================================================================
\textbf{질문(사슬 2번째 고리).} 다음으로 분극 $R_n$ 을 어떻게 독립 고정하는가? GITT(galvanostatic intermittent
titration)는 짧은 정전류 펄스 $+$ 긴 휴지의 반복으로 \emph{순간 응답}과 \emph{완화 응답}을 함께 준다 — 바로 이
시간척도 분리가 $R_n$ 과 $\chi_j$ 를 가르는 열쇠다. 다만 \textbf{GITT relaxation 을 곧바로 $k_j$ 의 직접 측정값으로
해석하지 않는다} — 여러 과정이 섞이기 때문이다:
\begin{equation}
\text{GITT relaxation}=\underbrace{\text{solid diffusion}}_{}+\underbrace{\text{electrolyte relaxation}}_{}+
\underbrace{\text{phase/staging relaxation}}_{\text{이론장 }k_j}+\underbrace{\text{OCV curvature}}_{}+\underbrace{\text{measurement filtering}}_{}.
\label{eq:ch6_gitt_decomp}
\end{equation}
\textbf{권장 절차(순차 분리) — 이걸로 무엇을 하나.}
\begin{enumerate}[label=G\arabic*),leftmargin=2.4em]
\item 펄스 직후 \emph{순간} step 으로 ohmic$\cdot$fast polarization 제약($R_n$ 의 일부; Ch3 $\eta_\ct$ 분리 시작).
\item 장시간 완화 끝점을 저속 OCV(\S\ref{sec:ch6_ocvfit})와 비교해 평형 골격 검증(O 제약과 정합 확인).
\item relaxation curve 를 단일 $k_j$ 로 단정하지 않고 single- vs multi-time-constant 후보를 비교한다(이론장
  spectrum 의 다중 $L$ — 여러 시간상수가 보이면 분포 $\rho_G$ 의 직접 증거).
\item C-rate series(\S\ref{sec:ch6_crate})와 결합해 $k_j$, diffusion limitation, transport limitation 을 분리.
\end{enumerate}
\textbf{$\chi_j$ vs $\eta_\ct$ 공선형(이론장 N 의 실현).} GITT 의 순간 step($\eta_\ct$, charge-transfer)과 느린
relaxation(barrier spectrum)을 분리한 \emph{후에만} $\chi_j(=\beta_j$; 조건부, 대칭 Marcus$\cdot$정상 target 극한 한정$)$
귀속이 성립한다. 즉 \textbf{$R_n$ 과 $\chi_j$ 는 GITT 의 시간척도 분리로 비로소 독립 제약}된다(\S\ref{sec:ch6_RnChi}
의 분리 전략을 데이터로 실현). 분극이 고정됐으니 \emph{다음 절}은 사슬 3번째 고리 — 다온도로 활성화 파라미터 제약 —
로 간다.

% ====================================================================
\section{S3 핵심: 다온도 Arrhenius 로 활성화 파라미터 제약}\label{sec:ch6_arrhenius}
% ====================================================================
\textbf{질문(사슬 3번째 고리).} 활성화 파라미터 $\Delta H_{a,j},\Delta S_{a,j},k_0$ 는 \emph{왜} 단일 온도로는
분리되지 않는가? $k_j=k_0 e^{-(\Delta H_a-T\Delta S_a-\chi_j\mathcal A_j)/RT}$ 에서 한 온도는 한 방정식만 주기
때문이다 — 여러 온도가 있어야 Arrhenius 직선의 \emph{기울기와 절편}이 따로 정해진다. 그래서 다온도 series 가 이
고리의 데이터다.

\textbf{Arrhenius 분해(무비약) — 무엇을 구하나.} 꼬리 길이 $L=|I|/(Q_\cell k_j)$ 의 로그를 취하면(이론장 L\_of\_G)
\begin{equation}
\ln\frac1L\;=\;\ln k_j-\ln\frac{|I|}{Q_\cell}\;=\;\ln k_0-\frac{\Delta H_{a,j}}{RT}+\frac{\Delta S_{a,j}}{R}
+\frac{\chi_j\mathcal A_j}{RT}-\ln\frac{|I|}{Q_\cell}.
\label{eq:ch6_arrhenius}
\end{equation}
\emph{한 줄씩}: $\Delta G_{a,j}=\Delta H_{a,j}-T\Delta S_{a,j}$(이론장 \AL{10})를 $k_j=k_0 e^{-(\Delta G_{a,j}-\chi_j
\mathcal A_j)/RT}$ 에 넣고 로그를 취했다. \textbf{$\mathcal A_j$ 를 상수로 고정하고}(같은 동작점) $\ln(1/L)$ 을 $1/T$ 에
대해 그리면, $\chi_j\mathcal A_j/(RT)=(\chi_j\mathcal A_j/R)\cdot(1/T)$ \emph{도} $1/T$ 에 선형이라 기울기에 들어간다.
따라서 기울기는 순수 $-\Delta H_{a,j}/R$ 가 \emph{아니라}
\begin{equation}
\frac{\dd\ln(1/L)}{\dd(1/T)}\Big|_{\mathcal A_j}=-\frac{\Delta H_{a,j}-\chi_j\mathcal A_j}{R}=-\frac{\Delta H_{a,j}^{\eff}}{R}
\label{eq:ch6_arrhenius_slope}
\end{equation}
즉 \emph{유효} 활성화 엔탈피 $\Delta H_{a,j}^{\eff}=\Delta H_{a,j}-\chi_j\mathcal A_j$ 를 준다.
\textbf{이걸로 무엇을 하나(순수 $\Delta H_a$ 분리).} \textbf{순수 $\Delta H_a$ 를 얻으려면} 구동력을 없앤 극한
$\mathcal A_j\to0$(꼬리 영역 $V_{n,\drive}\to U_j$, 무구동)로 외삽하거나, $\chi_j$ 가 독립 제약된 뒤 $\chi_j\mathcal A_j/R$
만큼 보정한다. $1/T\to0$ 외삽 절편 $=\ln k_0+\Delta S_{a,j}/R-\ln(|I|/Q_\cell)$ (prefactor$\cdot$엔트로피 묶음;
$\ln(|I|/Q_\cell)$ 상수항 포함)다. $k_0$ 와 $\Delta S_a$ 는 절편에서 공선형이므로 $k_0=(k_BT/h)\kappa$ 의
transition-state 절대값을 \emph{인용}(이론장 \AL{10}, \cite{eyring1935})해 $\Delta S_a$ 를 분리하되, $\kappa(T)$
미지로 $\ln\kappa$ 가 $\Delta S_a$ 로 흡수되는 잔여 공선형은 BOUNDED 다(\AL{69}; \cite{bardfaulkner}). $\Delta S_a$
를 reaction entropy(Ch2 \AL{24}, 별개 물리)와 혼동하지 않는다.
\begin{boundbox}
\textbf{$\mathcal A_j$ 고정의 의미(★ 정정).} 흔한 오해와 달리, $\chi_j\mathcal A_j/(RT)$ 항은 $\mathcal A_j$ 를
\emph{고정해도} $1/(RT)=(1/R)(1/T)$ 를 통해 기울기에 기여한다 — 따라서 고정 $\mathcal A_j$ Arrhenius 기울기는 순수
$-\Delta H_a/R$ 가 아니라 \emph{유효} $-(\Delta H_a-\chi_j\mathcal A_j)/R$ 다(식~\eqref{eq:ch6_arrhenius_slope}).
전위를 \emph{바꾸면} $\mathcal A_j$ 자체가 변해 기울기를 더 휘게 하므로, 깨끗한 분리는 한 동작점 기울기에서
$\chi_j\mathcal A_j/R$ 를 보정하거나 $\mathcal A_j\to0$ 외삽으로 한다. 이 $\chi_j\mathcal A_j$ 혼입이 N1
반증(\S\ref{sec:ch6_falsify})의 대상이다. (참고: $\mathcal A_j$ 를 일반형 $n_j^\eff F(V-U)=(RT/w_j)(V-U)$ 로 쓰면
$\chi_j\mathcal A_j/RT=\chi_j(V-U)/w_j$ 로 explicit $1/T$ 가 상쇄되나, 그땐 $w_j(T)$ 의 온도 표류가 같은 자리에 남아
— 어느 convention 이든 순수 $\Delta H_a$ 분리는 $\mathcal A_j\to0$/보정이 필요하다.)
\end{boundbox}
\textbf{Arrhenius 직선성 = 반증 가능성(이론장 N2).} $\ln(1/L)$ 대 $1/T$ 가 \emph{직선이 아니면} 단일 활성화 그림이
기각된다 — 분포 $\rho_G$ 의 폭이 온도에 따라 유효 기울기를 휘게 하거나(stretched, \AL{14}), transport 가
지배하는 경우다. 직선성 자체가 모델의 검정이다 — 이 검정의 데이터 실현이 곧 \S\ref{sec:ch6_falsify} 의 N1 이다.
활성화 파라미터가 고정됐으니 \emph{다음 절}은 사슬 마지막 고리 — C-rate series 로 전위 의존 가속과 꼬리 제약 — 이다.

% ====================================================================
\section{S3--S5 검증: C-rate series 로 전위 의존 가속과 꼬리 제약}\label{sec:ch6_crate}
% ====================================================================
\textbf{질문(사슬 4번째 고리).} 마지막으로 $\chi_j$ 와 병목 $k_{j,\lim}$ 을 무엇으로 닫는가? $0.1/0.2\,\mathrm C$ 는
quasi-equilibrium proxy(peak 위치/폭), $1.0\,\mathrm C$ 는 kinetic-lag/transport stress test 다(strict frozen limit
아님). \textbf{검증 순서(순차 제약의 마지막 단계).} 각 단계가 앞 단계에서 고정한 파라미터 위에서 새 파라미터만
자유화한다:
\begin{enumerate}[label=C\arabic*),leftmargin=2.4em]
\item 저속 OCV 로 평형 골격($U_j,w_j,Q_{j,\tot},Q_\bg$) 제약(\S\ref{sec:ch6_ocvfit}).
\item GITT 로 fast polarization($R_n$)$\cdot$relaxation signature 제약(\S\ref{sec:ch6_gittfit}).
\item 다온도로 활성화 파라미터($\Delta H_a,\Delta S_a,k_0$) 제약(\S\ref{sec:ch6_arrhenius}).
\item $0.1/0.2\,\mathrm C$ 에서 ICA/DVA peak 위치$\cdot$폭 검증(S1--S2; 평형 골격 재확인).
\item $0.5/1.0\,\mathrm C$ 에서 \textbf{꼬리$\cdot$broadening$\cdot$shift} 로 $\chi_j$ 제약(S3--S6, kernel integral).
  전위 의존 $\partial\ln L/\partial V_{n,\drive}<0$ 가 C-rate 증가 시 꼬리 단축으로 나타나는지 확인(\AL{62}).
\item 잔차를 $\rho_G(g),k_{j,\lim},R_n$, transport 중 \emph{어느 항으로} 설명 가능한지 독립 제약과 함께 판단
  (한 항으로 모두 흡수 금지; 식~\eqref{eq:ch6_no_free}).
\end{enumerate}
\textbf{이 순서가 식별성을 보장하는 이유(의미).} 각 단계가 \emph{앞 단계에서 고정된 파라미터} 위에서 \emph{새 파라미터
하나(군)} 만 자유화하므로, 식~\eqref{eq:ch6_no_free} 의 공선형 집합이 한 번에 풀리지 않고 차례로 닫힌다 — 이것이
``동시 자유 피팅'' 과 ``순차 식별'' 의 결정적 차이다. 순차 식별이 끝났으니, \emph{다음 절}은 ``맞췄다''로 끝나지
않게 하는 마지막 단계 — 반증 — 으로 간다.

% ====================================================================
\section{반증(falsification): forward-only 예측의 데이터 검정}\label{sec:ch6_falsify}
% ====================================================================
\textbf{질문.} 순차 식별로 곡선을 맞췄다 치자 — 그런데 ``곡선을 맞췄다''가 곧 ``모델이 옳다''인가? 아니다. 충분한
자유도면 \emph{어떤} 모델도 맞출 수 있어 의미가 없다. \textbf{모델의 가치는 \emph{반증 가능성}}이다. 이론장 Ch1--5
는 모두 \emph{forward-only}(파라미터$\to$예측) 모델이므로, 본 절은 이론장의 반증 규칙 N1--N4 를 \emph{데이터 시험}으로
실현한다 — 즉 ``무엇이 어긋나면 이 귀속을 \emph{기각}하는가''를 못박는다.
\begin{loopbox}
\textbf{도입 관찰의 반증 가능 시험(이론장 N1--N4 의 데이터 실현).} 도입 관찰 ``저온$\cdot$고율서 긴 겹침 꼬리''를
다음으로 시험한다(어느 하나라도 어긋나면 해당 귀속을 \emph{기각}).
\begin{itemize}[leftmargin=1.6em]
\item \textbf{(N1) Arrhenius 기울기 vs $E_D$.} 꼬리 길이의 $\ln(1/L)$ 대 $1/T$ 기울기(식~\eqref{eq:ch6_arrhenius})가
  독립 측정 확산 활성화 $E_D$ 와 \emph{일치}하고 \emph{전위 의존이 없으면}, 꼬리는 barrier-spectrum 이 아니라
  transport 지배다 $\Rightarrow$ barrier-spectrum 귀속 \emph{기각}. C-rate$\times$온도 교차표로 측정.
\item \textbf{(N2) 전위 의존 spectrum shift.} 전위/$|I|$ 증가 시 spectrum 이 short-$L$ 로 이동하는지
  ($\partial\ln L/\partial V_{n,\drive}<0$, 이론장 L\_of\_G). 이동이 없으면 $\chi_j$ 항 \emph{기각}(가속의 전위
  의존이 모델의 핵심 예측).
\item \textbf{(N3) 저온 lineshape.} 저속 OCV near-peak lineshape 이 저온에서 \emph{안 좁아지면}, 평형 broadening
  ($w=RT/F$ 는 저온서 좁아짐, 이론장 dxidV)으로 설명 안 되는 heterogeneous disorder 기여가 존재 $\Rightarrow$
  kinetic 단독 귀속 기각, disorder 와 분리(이론장 flagbox).
\item \textbf{(N4) $\rho_G$ forward-only.} $\rho_G(g)$ 를 관측 꼬리에서 \emph{역산하지 않고}(비유일, \AL{16})
  독립 근거(입자 크기 분포$\cdot$EIS$\cdot$계산)로 주어 forward 예측만 한다 — 역산하면 비유일이라 반증력이 없다.
\end{itemize}
이로써 \textbf{도입 관찰 $\to$ 인과 사슬(이론장 Ch1) $\to$ 발열/속도론/히스테리시스(Ch2--5) $\to$ 피팅식 $\to$ 반증
가능한 데이터 시험(본 부록)} 으로 한 바퀴가 닫힌다.
\end{loopbox}
\begin{boundbox}
\textbf{$\rho_G$ 역산 금지의 피팅 함의(\AL{16}).} 피팅 도구는 잔차를 줄이려 $\rho_G$ 를 자유화하고 싶어 하지만,
stretched 거동에서 $\rho_G$ 로의 역매핑은 비유일(재포획 등 다른 메커니즘도 같은 꼬리)이다. 따라서 $\rho_G$ 를
피팅 변수로 풀면 반증력이 사라진다 — 독립 입력으로 고정하고, 그 입력이 데이터를 \emph{예측}하는지만 본다.
\end{boundbox}

반증 규칙이 섰으니, \emph{다음 절}은 반증되지 않은 항이라도 \emph{데이터가 침묵하는} 발열 항을 어떻게 다룰지를 본다.

% ====================================================================
\section{calorimetry 부재 시 발열 항의 forward-prediction 제한}\label{sec:ch6_heat}
% ====================================================================
\textbf{질문.} 발열 항은 왜 피팅하지 않는가? \textbf{이론장 Ch2/Ch4 의 발열 항을 ICA/전압 데이터만으로 피팅하지
않는다}(\AL{67}). 비가역 발열 $\dot{\mathcal Q}_\irr=\sum_j n_j^{\eff}I_j\eta_j\ge0$(Ch4)과 가역 발열
$\dot{\mathcal Q}_\rev$(Ch2 Bernardi energy balance \cite{bernardi1985}; 다공성 전극 열모델
\cite{thomasnewman2003})은 \emph{전류$\cdot$과전압$\cdot$entropy coefficient} 로부터 \textbf{forward 로 예측}만
하고, 그 예측을 \emph{calorimetry} 가 확보됐을 때 검증한다(\AL{67}).
\textbf{이유(의미).} ICA/전압 데이터는 발열을 직접 보지 않으므로, 발열 파라미터를 전압 잔차에 피팅하면 다른 전압
효과(분극 등)와 공선형이 되어 가짜 값을 얻는다(식~\eqref{eq:ch6_no_free} 의 정신을 발열로 확장) — 즉 ``보지 않는
양을 피팅하지 않는다''는 식별성 규칙의 발열판이다.
\begin{boundbox}
\textbf{발열의 ICA 되먹임(등온 한계의 신호).} 등온 가정이 깨지면(고율 자기발열) 이론장 Ch.2 의 되먹임
($\dot{\mathcal Q}\to T\uparrow\to k_j\uparrow\to L\downarrow$ 꼬리 단축, 동시에 $w=RT/F\uparrow$ 로 평형 peak
폭 증가)이 ICA tail 에 두 부호로 작용한다. 따라서 고율 ICA tail 을 단일 부호로 단정하지 않고, 등온이 의심되면
Ch2/Ch4 thermal coupling 으로 돌아간다(\AL{67}).
\end{boundbox}

음극 방전 피팅의 전 과정이 섰으니, \emph{다음 절}은 이를 충전$\cdot$양극$\cdot$full-cell 로 확장하는 조건을 본다.

% ====================================================================
\section{충전$\cdot$양극$\cdot$full-cell 로의 확장 조건 (음극 방전 검증 후)}\label{sec:ch6_extension}
% ====================================================================
\textbf{질문.} 음극 방전 모델을 언제, 어떤 순서로 확장하는가? 핵심은 \emph{음극 방전이 먼저 식별$\cdot$검증된
뒤에만} 확장한다는 것이다.
\begin{enumerate}[label=E\arabic*),leftmargin=2.4em]
\item \textbf{충전 branch.} 이론장 Ch.5 의 signed current $I_s$ 와 metastable branch target
  $\xi_{\tar,j}^b=\xi_{\sstat,j}^b$, 유도된 부호 $s_{\phi,j}^b=-1$ 을 사용한다(\AL{53}--56). branch switching 시
  $\xi_j,h_j,V_n$ continuity(재초기화 X), branch 별 root finding(\S\ref{sec:ch6_S1})으로 매 시점 $V_n$ 결정.
\item \textbf{히스테리시스 분리 검증.} $\Delta V_\obs\ne\Delta V_\hys$(Ch5 \AL{58}): 관측 loop 폭에서 polarization
  ($R_n^b$; $|I_s|\to0$ 외삽으로 소멸)$\cdot$kinetic lag$\cdot$transport 를 빼고 남는 잔존 분지 차만 true
  hysteresis 후보로 둔다. empirical branch fit 은 metastable free-energy/domain memory 분리 안 되면 \emph{validation
  candidate} 로만 취급(주 결론서 제외, Ch5 forward-only).
\item \textbf{양극.} 양극에 별도 전하 보존식$\cdot$STO--OCV$\cdot$GITT 를 같은 절차로 적용.
\item \textbf{full-cell.} $V_\cell=V_p-V_n-\eta_\loss$ 또는 apparent electrode voltage 표현; loss/polarization
  이중 계상 금지(Ch5). full-cell ICA/DVA 는 양극/음극/branch signed loss/transport/polarization 이 섞인 관측
  신호라 음극 단독과 직접 같지 않다.
\item \textbf{순서.} \textbf{음극 방전이 먼저 식별$\cdot$검증되지 않으면 full-cell 조합 해석으로 넘어가지 않는다.}
\end{enumerate}
\textbf{의미.} 이 확장 게이트는 ``검증된 단순 단위(음극 방전)부터 합성한다''는 식별성 원칙의 시스템 수준 적용이다.
\emph{다음 절}은 이 전 과정에서 \emph{쓰면 안 되는} 피팅 편의를 한 자리에 격리한다.

% ====================================================================
\section{금지되는 피팅 편의와 실무 보조의 격리}\label{sec:ch6_forbidden}
% ====================================================================
\textbf{질문.} 무엇을 본문 물리식으로 \emph{쓰면 안 되는가}? 잔차를 줄이려는 수치 편의가 모델 물리식으로 둔갑하면
식별성과 반증력이 동시에 무너진다. 다음은 금지 목록이다(본문 물리식으로 쓰지 않음, GS-4).
\begin{enumerate}[label=(\arabic*),leftmargin=2.4em]
\item 강구동 rate 폭주를 물리 의미 없는 cap/clip 으로 자르기(대신 stiff solver $+$ 물리 병목 $k_{j,\lim}$).
\item 잔차를 arbitrary background correction(임의 wiggle)으로 모두 흡수(O6).
\item branch hysteresis 를 단순 charge/discharge offset 으로만 피팅(Ch5 metastable 분리 없이).
\item negative quadrature weight/비물리 time constant 로 relaxation spectrum 맞추기(완전 단조성 위반, \S\ref{sec:ch6_reduction}).
\item Fredholm/Padé(Plan A) correction factor 를 불분명한 fitting knob 으로 사용(\S\ref{sec:ch6_closure}).
\item $\rho_G$ 를 관측 꼬리에서 역산(비유일, N4).
\item 근거 없는 수치 임계값($\varepsilon_\tol$ 등)을 established 로 사용(\S\ref{sec:ch6_epsilon} 의 측정 기반 정의로 대체).
\end{enumerate}
\textbf{의미.} 이 목록은 앞 절들의 모든 금지 조항을 한 곳에 모은 ``음성 목록''이며, 자체 검수표(\S\ref{sec:ch6_selfcheck})가
이를 점검한다. 단, 디버깅용 수치 보조까지 금지하는 것은 아니다 — box 가 그 격리 경계를 정한다.
\begin{practicebox}
\textbf{실무 보조(본문 물리식 아님, GS-4$''$).} 초기값 탐색$\cdot$디버깅에서 bounded transform, 임시 clipped rate,
smoothed background 같은 수치 보조를 \emph{임시} 사용할 수 있으나 본 모델 물리식이 아니다. 최종 해석$\cdot$논문용
식$\cdot$검증 결과에서는 제거하고 물리적 제한항(병목 $k_{j,\lim}$) 또는 검증된 축약법으로 대체한다. 피팅 보조가
결과에 남았는지 여부는 \S\ref{sec:ch6_selfcheck} 검수표가 점검한다.
\end{practicebox}

% ====================================================================
\section{부록 B 자체 검수표 (구 Chapter 6)}\label{sec:ch6_selfcheck}
% ====================================================================
{\small
\begin{longtable}{p{0.74\textwidth} p{0.16\textwidth}}
\toprule 검수 항목 & 결과 \\ \midrule \endhead
이론장 Ch1--5 본문 수정 없이 부록만 적층 & 통과(\S\ref{sec:ch6_inherit}) \\
★ 피팅 실무가 주역, 수치해법은 inner-loop solver 로 종속(사용자 5-30) & 통과(\S\ref{sec:ch6_S},\ref{sec:ch6_dae}) \\
피팅 평가 순서 S1--S6 을 피팅 워크플로로 실현 & 통과(\S\ref{sec:ch6_S}) \\
피팅 파라미터 전수 목록 $+$ 정의 장$\cdot$1차 제약 실험 매핑 & 통과(\S\ref{sec:ch6_inherit},\ref{sec:ch6_identif}) \\
★ 식별성: $R_n$ vs $\chi_j$ 공선형 분리(독립 실험 순차) & 통과(\S\ref{sec:ch6_RnChi}) \\
★ $R_n\ne R_\ct\ne R_{n,\eff}$ 분리 유지(이론장 Ch1$\cdot$3$\cdot$4 계승) & 통과(\S\ref{sec:ch6_threeR} keybox) \\
동시 자유화 금지 집합 $+$ 순차 제약 사슬(OCV$\to$GITT$\to$Arrhenius$\to$C-rate) & 통과(\S\ref{sec:ch6_noFree}) \\
dynamic root vs OCV root 의 derivative 구분(식~\eqref{eq:ch6_dyn_deriv} vs \eqref{eq:ch6_ocv_root}) & 통과(\S\ref{sec:ch6_S1}) \\
OCV root 분모 $=$ 이론장 Ch2 평형 chemical capacitance 정합 & 통과(\S\ref{sec:ch6_S1}) \\
전하 보존식 해 존재 조건을 residual 로 흡수 X(비물리 파라미터 reject) & 통과(\S\ref{sec:ch6_S1},\ref{sec:ch6_dae}) \\
index-1 DAE \emph{왜 index-1 인지} 유도(단조성 floor) & 통과(\S\ref{sec:ch6_dae}) \\
강구동 rate 를 cap 으로 자르지 않고 stiff$+$물리 병목으로 처리 & 통과(\S\ref{sec:ch6_dae} boundbox) \\
직접 $g$-grid 를 기준 해법으로; 모든 축약은 대비 검증 & 통과(\S\ref{sec:ch6_grid}) \\
이중 계상(목표 vs 잔차) 수치 점검 극한 & 통과(\S\ref{sec:ch6_grid} verifybox) \\
Prony/rational $a_\ell\ge0,\tau_\ell>0$ 을 완전 단조성으로 유도(편의 아님) & 통과(\S\ref{sec:ch6_reduction}) \\
★ FLAGGED 허위 정정: $\varepsilon_\tol$ 을 측정 기반(이산화 오차$\vee$잡음)으로 재정의 & 통과(\S\ref{sec:ch6_epsilon} eq~\eqref{eq:ch6_epsilon}) \\
★ $\varepsilon_\tol$ 특정 숫자값을 flagbox 로 정직 표기(established 사용 X) & 통과(\S\ref{sec:ch6_epsilon} flagbox) \\
★ root-find/solver tolerance 도 물리 scale$\cdot$잡음으로 근거화(GS-4$'$) & 통과(verifybox \S\ref{sec:ch6_S1},\ref{sec:ch6_dae}) \\
Plan A(Fredholm)를 M1--M5$\cdot\varepsilon$ 가속 후보로 제한; broad-kernel/Volterra 한계 명시 & 통과(\S\ref{sec:ch6_closure}) \\
저속 OCV 골격(O1--O6); $w_j$ vs $\rho_G(g)$ 구분, $n^{\eff}w_j=RT/F$ 정합 & 통과(\S\ref{sec:ch6_ocvfit}) \\
GITT relaxation 을 $k_j$ 직접 측정으로 단정 X; $\chi$ vs $\eta_\ct$ 분리 & 통과(\S\ref{sec:ch6_gittfit}) \\
★ 다온도 Arrhenius 로 $\Delta H_a$ 기울기 식별; $\mathcal A_j$ 고정 조건 명시 & 통과(\S\ref{sec:ch6_arrhenius}) \\
$0.1\mathrm C/1.0\mathrm C$ 를 strict limit 아닌 proxy/stress test 로 & 통과(\S\ref{sec:ch6_crate}) \\
★ 도입 관찰 $\to$ 반증 N1--N4 데이터 시험으로 loop 폐합 & 통과(\S\ref{sec:ch6_falsify} loopbox) \\
calorimetry 없이 발열 파라미터 피팅 X(forward prediction) & 통과(\S\ref{sec:ch6_heat}) \\
피팅 편의(cap/clip/transform)를 본문서 배제, practicebox 격리 & 통과(\S\ref{sec:ch6_S6},\ref{sec:ch6_forbidden}) \\
충전/양극/full-cell 확장을 음극 방전 검증 이후로 보류 & 통과(\S\ref{sec:ch6_extension}) \\
모든 가정 AL-60$\sim$69 인용; 무근거 established 0; FLAGGED 정직 표기 & 통과(\S\ref{sec:ch6_al}) \\
``자명/clearly/쉽게/obviously'' 본문 0건(G-noleap) & 통과 \\
★ clean-spine: box 제외 본문 줄기 연속(G-spine)$\cdot$동기(G-motive)$\cdot$의미(G-interpret)$\cdot$다리(G-bridge) & 통과(전 절 5비트) \\
\bottomrule
\end{longtable}
}

% ====================================================================
\section{부록 B 가정 원장 (구 Chapter 6, Assumption Ledger, AL-60$\sim$69)}\label{sec:ch6_al}
% ====================================================================
본 부록이 사용한 가정과 그 tier$\cdot$근거$\cdot$검증 DOI 다(통합 master 와 정합; AL-64$\sim$69 가 본 부록 신규 등재).
{\small
\begin{longtable}{p{0.05\textwidth} p{0.49\textwidth} p{0.12\textwidth} p{0.26\textwidth}}
\toprule AL & 가정 진술 & tier & 근거(검증 DOI) \\ \midrule \endhead
60 & index-1 DAE(root-constrained ODE), BDF/Radau 시간 적분; 단조성 floor 가 index-1 보장 & GROUNDED & brenan1996(10.1137/\brk 1.9781611971224), hindmarsh2005(10.1145/\brk 1089014.1089020) \\
61 & Plan A closure: Fredholm 2종 ratio-substitution closed-form(사용자 PhD) & BOUNDED & jcp2017(10.1063/\brk 1.5000882), lee2011(10.1063/\brk 1.3565476), son2013(10.1063/\brk 1.4802584) \\
62 & ★ ratio ansatz 의 broad-kernel 열화 $=$ stretched-tail(저온) 영역서 부정확 $\to$ 직접 grid validator 필수 & BOUNDED & jcp2017 Eq.34 자기명시(10.1063/\brk 1.5000882) \\
63 & ICA($dQ/dV$) rate-dependency$\cdot$미분 잡음 증폭$\cdot$best-practice & GROUNDED & dubarry2022(10.3389/\brk fenrg.2022.1023555), fly2020(10.1016/\brk j.est.2020.101329) \\
64 & ★(신규) root-find/solver 종료 tolerance 는 물리 scale($\delta_q Q_\cell$, $\xi\in[0,1]$)$\cdot$기계 epsilon floor 로 정함(임의 상수 아님) & BOUNDED & brenan1996(10.1137/\brk 1.9781611971224); 수치해석 표준 \\
65 & ★(신규) 분포 적분 축약(adaptive quadrature/moment/Prony)은 기준 grid 대비 검증 시만; Prony $a_\ell\ge0,\tau_\ell>0$ 은 완전 단조성 보존 & BOUNDED & lindsey1980(10.1063/\brk 1.440530, KWW$\leftrightarrow$분포), johnston2006(10.1103/\brk PhysRevB.74.184430, 연속합) \\
66 & ★(신규) closure 검증 tolerance $\varepsilon_\tol=\max(\varepsilon_{\mathrm{disc}},\varepsilon_\noise)$; \emph{특정 숫자값은 데이터 의존 미확정} & FLAGGED(값) / BOUNDED(정의) & dubarry2022(10.3389/\brk fenrg.2022.1023555), fly2020(10.1016/\brk j.est.2020.101329); 값은 사용자 데이터서 산출 \\
67 & ★(신규) calorimetry 부재 시 Ch2/Ch4 발열 항은 forward prediction 으로만(전압 데이터에 발열 파라미터 피팅 X) & BOUNDED & bernardi1985(10.1149/\brk 1.2113792), thomasnewman2003(10.1149/\brk 1.1531194) \\
68 & ★(신규) 식별성: $\{R_n,k_j^{\eff},w_j,\rho_G,Q_\bg,Q_{j,\tot},k_{j,\lim}\}$ 동시 자유화 금지; 순차 제약(OCV$\to$GITT$\to$Arrhenius$\to$C-rate)으로만 식별 & BOUNDED & bardfaulkner(ISBN 978-0-\brk 471-04372-0), fly2020(10.1016/\brk j.est.2020.101329) \\
69 & ★(신규) Arrhenius 분해: $\mathcal A_j$ 고정 기울기 $=-(\Delta H_a-\chi_j\mathcal A_j)/R$ = \emph{유효} 엔탈피($\chi\mathcal A/RT$ 가 $1/T$-선형 기여); 순수 $\Delta H_a$ 는 $\mathcal A_j\to0$ 외삽/보정; $k_0,\Delta S_a$ 절편 공선형($\kappa$ 미지 잔여) & BOUNDED & eyring1935(10.1063/\brk 1.1749604), bardfaulkner(ISBN 978-0-\brk 471-04372-0) \\
\bottomrule
\end{longtable}
}
\textbf{tier 요약.} GROUNDED: AL-60,63,69(분해). BOUNDED: AL-61,62,64,65,67,68,69(분리). \textbf{FLAGGED: AL-66 의
$\varepsilon_\tol$ \emph{특정 숫자값}}(정의는 측정 기반으로 BOUNDED, 값만 데이터 확보 후 확정). 본 부록은 무근거
established 수치 임계값을 본문에 두지 않는다(구 consolidated 판의 $\varepsilon_\tol$ 허위를 \S\ref{sec:ch6_epsilon}
에서 정정).

% ====================================================================
\section{부록 B 의 결론 (구 Chapter 6)}\label{sec:ch6_concl}
% ====================================================================
첫째, \textbf{본 부록(구 Chapter 6)은 이론장 유도식(Ch1--5)을 ICA/DVA 데이터에 \emph{피팅하는 방법론}이다}. 피팅 평가
순서 S1--S6 (\S\ref{sec:ch6_S})이 한 점의 $\dd Q/\dd V_n$ 을 주는 forward 계산이고, 이를 측정 ICA 와 비교해
파라미터를 식별한다. \textbf{수치 해법(전하 보존식 root-find, index-1 DAE 시간 적분, self-consistent Volterra 의
직접 grid 풀이)은 이 피팅의 inner-loop solver 로 \emph{종속}}되며 주역이 아니다(사용자 5-30 의도 반영).

둘째, \textbf{피팅의 본질은 회귀가 아니라 순차 식별}이다(\S\ref{sec:ch6_identif}). $R_n$ 과 $\chi_j$ 처럼 같은
꼬리를 만드는 공선형 파라미터는 한 데이터에서 분리되지 않으므로, 저속 OCV$\to$GITT$\to$다온도 Arrhenius$\to$
C-rate series 를 순서대로 걸어 식~\eqref{eq:ch6_no_free} 의 동시 자유화 금지 집합을 차례로 닫는다. $R_n\ne R_\ct
\ne R_{n,\eff}$ 의 세 저항 분리(이론장 Ch1$\cdot$3$\cdot$4 계승)가 식별성과 발열 예측을 동시에 지킨다.

셋째, \textbf{전하 보존식 root 의 dynamic/OCV 두 형태는 derivative 가 다르므로 분리}하며(식~\eqref{eq:ch6_dyn_deriv}
vs \eqref{eq:ch6_ocv_root}), 후자의 분모는 이론장 Ch.2 평형 chemical capacitance 와 정합한다. index-1 DAE 는 단조성
floor 가 index 를 1 로 만든다는 \emph{유도}로 정당화되며, 강구동 stiffness 는 cap 이 아니라 implicit solver 와
물리 병목으로 처리한다.

넷째, \textbf{closure 의 기준 해법은 직접 $g$-grid}(\S\ref{sec:ch6_grid})이고, adaptive quadrature/moment matching/
Prony-rational/Plan A(Fredholm)는 모두 기준 해 대비 검증될 때만 쓴다. Plan A 는 문헌적 가속 후보이되 흑연의 넓은
spectrum(저온 stretched tail)에서 가장 부정확할 수 있어(\AL{62}, jcp2017 자기명시) 항상 직접 grid 가 validator 다.

다섯째, \textbf{★ 구 consolidated 판의 $\varepsilon_\tol$ 허위(근거 없는 수렴 임계값, flagbox 0건)를 정정}했다
(\S\ref{sec:ch6_epsilon}). $\varepsilon_\tol$ 을 ``작으면 좋다''가 아니라 기준 해의 이산화 오차와 측정 잡음으로
인한 ICA 불확실도의 \emph{큰 쪽}(식~\eqref{eq:ch6_epsilon})으로 \emph{측정 기반} 재정의하고, 그 \emph{특정 숫자값}은
데이터 의존이라 FLAGGED 로 정직 표기했다. root-find$\cdot$solver tolerance 도 같은 원리로 물리 scale$\cdot$잡음
floor 에서 근거화했다(\AL{64}, GS-4$'$).

여섯째, \textbf{반증(forward-only)이 모델의 가치}다(\S\ref{sec:ch6_falsify}). Arrhenius 직선성(N1), 전위 의존
spectrum shift(N2), 저온 lineshape(N3), $\rho_G$ 역산 금지(N4)로 도입 관찰을 \emph{기각 가능}하게 시험한다.
calorimetry 부재 시 발열 항은 피팅하지 않고 forward prediction 으로만 둔다(\AL{67}).

일곱째, \textbf{도입 관찰(저온$\cdot$고율 긴 겹침 꼬리)} $\to$ 인과 사슬(이론장 Ch1) $\to$ 발열/속도론/히스테리시스
(Ch2--5) $\to$ 피팅식과 순차 식별$\to$ 반증 가능 데이터 시험(본 부록 N1--N4)으로 \emph{한 바퀴가 닫힌다}. 음극
방전 모델이 먼저 식별$\cdot$검증된 뒤에야 충전 branch$\cdot$양극$\cdot$full-cell 로 확장한다.

% ====================================================================
% 참고문헌 (본 부록이 인용하는 항목만; 이론장 master bib 와 정합)
% ====================================================================
\begin{thebibliography}{99}
\bibitem{eyring1935} H. Eyring, ``The Activated Complex in Chemical Reactions,'' \emph{J. Chem. Phys.}
  \textbf{3}, 107 (1935). DOI: 10.1063/1.1749604.
\bibitem{bardfaulkner} A. J. Bard, L. R. Faulkner, \emph{Electrochemical Methods}, 2nd ed., Wiley (2001).
  ISBN 978-0-471-04372-0.
\bibitem{johnston2006} D. C. Johnston, ``Stretched Exponential Relaxation Arising from a Continuous Sum
  of Exponential Decays,'' \emph{Phys. Rev. B} \textbf{74}, 184430 (2006). DOI: 10.1103/PhysRevB.74.184430.
\bibitem{lindsey1980} C. P. Lindsey, G. D. Patterson, ``Detailed Comparison of the Williams-Watts and
  Cole-Davidson Functions,'' \emph{J. Chem. Phys.} \textbf{73}, 3348 (1980). DOI: 10.1063/1.440530.
\bibitem{jcp2017} (사용자 PhD) K. Lee, S. Lee, C. H. Choi, S. Lee, ``Effects of External Electric Field and Anisotropic Long-Range Reactivity on Charge Separation Probability,''
  \emph{J. Chem. Phys.} \textbf{147}, 144111 (2017). DOI: 10.1063/1.5000882.
\bibitem{lee2011} (Refs.~6) S. Lee, C. Y. Son, J. Sung, S.-H. Chong, ``Communication: Propagator for Diffusive Dynamics of an Interacting Molecular Pair,''
  \emph{J. Chem. Phys.} \textbf{134}, 121102 (2011). DOI: 10.1063/1.3565476.
\bibitem{son2013} (Refs.~7) C. Y. Son, J. Kim, J.-H. Kim, J. S. Kim, S. Lee, ``An Accurate Expression for the Rates of Diffusion-Influenced Bimolecular Reactions with Long-Range Reactivity,''
  \emph{J. Chem. Phys.} \textbf{138}, 164123 (2013). DOI: 10.1063/1.4802584.
\bibitem{dubarry2022} M. Dubarry, D. Anseán, ``Best practices for incremental capacity analysis,''
  \emph{Front. Energy Res.} \textbf{10}, 1023555 (2022). DOI: 10.3389/fenrg.2022.1023555.
\bibitem{fly2020} A. Fly, R. Chen, ``Rate dependency of incremental capacity analysis (dQ/dV)...,''
  \emph{J. Energy Storage} \textbf{29}, 101329 (2020). DOI: 10.1016/j.est.2020.101329.
\bibitem{brenan1996} K. E. Brenan, S. L. Campbell, L. R. Petzold, \emph{Numerical Solution of Initial-Value
  Problems in Differential-Algebraic Equations}, SIAM Classics in Applied Mathematics 14 (1996).
  DOI: 10.1137/\brk 1.9781611971224.
\bibitem{hindmarsh2005} A. C. Hindmarsh, P. N. Brown, K. E. Grant, S. L. Lee, R. Serban, D. E. Shumaker,
  C. S. Woodward, ``SUNDIALS: Suite of Nonlinear and Differential/Algebraic Equation Solvers,''
  \emph{ACM Trans. Math. Softw.} \textbf{31}, 363 (2005). DOI: 10.1145/\brk 1089014.1089020.
\bibitem{bernardi1985} D. Bernardi, E. Pawlikowski, J. Newman, ``A General Energy Balance for Battery
  Systems,'' \emph{J. Electrochem. Soc.} \textbf{132}, 5 (1985). DOI: 10.1149/\brk 1.2113792.
\bibitem{thomasnewman2003} K. E. Thomas, J. Newman, ``Thermal Modeling of Porous Insertion Electrodes,''
  \emph{J. Electrochem. Soc.} \textbf{150}, A176 (2003). DOI: 10.1149/\brk 1.1531194.
\end{thebibliography}

\end{document}
